% Research continuation for the ProveIt repository, 11 October 2026.
% This file is assembled with the reviewed companion sections into the
% standalone deliverable Higher_Digamma_Powers_and_Ordered_Gamma_Sums.tex.
\documentclass[11pt,a4paper]{article}
\usepackage[T1]{fontenc}
\usepackage[utf8]{inputenc}
\usepackage{lmodern}
\usepackage[margin=27mm,headheight=15pt]{geometry}
\usepackage{amsmath,amssymb,amsthm,mathtools,mathrsfs}
\usepackage{booktabs,longtable,array}
\usepackage{microtype,xcolor,enumitem,fancyhdr}
\usepackage{xurl}
\usepackage[colorlinks=true,linkcolor=blue!45!black,citecolor=blue!45!black,urlcolor=blue!45!black]{hyperref}
\hypersetup{pdftitle={Higher Digamma Powers and Ordered Gamma Sums},
 pdfauthor={Research continuation prepared with ChatGPT for Vladimir Reshetnikov and ProveIt},
 pdfsubject={Exact finite parts, spectral zeta coefficients, and polylogarithmic depth cancellations}}
\numberwithin{equation}{section}
\newtheorem{theorem}{Theorem}[section]
\newtheorem{proposition}[theorem]{Proposition}
\newtheorem{lemma}[theorem]{Lemma}
\newtheorem{corollary}[theorem]{Corollary}
\theoremstyle{definition}
\newtheorem{definition}[theorem]{Definition}
\theoremstyle{remark}
\newtheorem{remark}[theorem]{Remark}
\DeclareMathOperator{\FP}{FP}
\DeclareMathOperator{\CT}{CT}
\DeclareMathOperator*{\Res}{Res}
\DeclareMathOperator{\Li}{Li}
\newcommand{\dd}{\mathrm d}
\newcommand{\Hrel}{\mathcal H}
\newcommand{\C}{\mathbb C}
\newcommand{\N}{\mathbb N}
\newcommand{\Q}{\mathbb Q}
\newcommand{\Z}{\mathbb Z}
\newcommand{\Log}{\operatorname{Log}}
\newcommand{\zet}{\zeta}
\newcommand{\commit}{b7667ae40a2957eed9a53d8559d53a3ee20dd1c1}
\setlength{\parskip}{0.35em}
\setlength{\emergencystretch}{2em}
\setlist{itemsep=3pt,topsep=5pt}
\allowdisplaybreaks[2]
\pagestyle{fancy}
\fancyhf{}
\fancyhead[L]{\small Higher Digamma Powers and Ordered Gamma Sums}
\fancyhead[R]{\small ProveIt research continuation}
\fancyfoot[C]{\thepage}

\begin{document}
\hypersetup{pageanchor=false}
\begin{titlepage}
\vspace*{14mm}
\noindent{\large\sffamily PROVEIT\quad /\quad RESEARCH CONTINUATION}
\vspace{18mm}

\noindent{\Huge\bfseries Higher Digamma Powers\\[3mm]
and Ordered Gamma Sums}
\vspace{8mm}

\noindent{\Large Exact finite parts, spectral zeta coefficients,\\[2mm]
and polylogarithmic depth cancellations}
\vspace{15mm}

\noindent Prepared with ChatGPT for Vladimir Reshetnikov and the ProveIt project\\
11 October 2026
\vspace{12mm}

\begin{abstract}
We develop a complete finite-part formula for every positive integer power
of the digamma function.  A normalized Mittag--Leffler expansion, finite
summation by parts, and a large-circle residue calculation reduce the
unit-interval moment to an absolutely convergent harmonic sum and the
linear Laurent coefficient of an explicit harmonic Dirichlet series.
The entire boundary polynomial is evaluated; its constant is zero for
odd powers and $(-1)^m\pi^{2m}/(2m+1)$ for power $2m$.  A second form
uses only an ordinary convergent series and an explicit finite list of
ordinary Stieltjes constants.  We carry out the fifth power completely,
including the mixed double-zeta derivatives that do not occur at powers
two through four, and give normalized primitives and differentiated
polygamma identities for all powers.

A complementary two-marker Gamma formula sums ordered relative harmonic
words over depth while retaining an arbitrary complex tagged exponent.
Integer marker differences give finite Hurwitz-zeta reductions, and
spectral differentiation produces generalized Stieltjes and log-Gamma
identities.  We also prove finite-interval formulas with several tagged
exponents and polylogarithmic depth cancellations.  Two precisely located
external-source errors are corrected by direct proofs.  Exact algebraic
checks and independent high-precision diagnostics accompany the proofs.
The Gaussian $S_6$ and revised $S_8$ conjectures remain unresolved.
\end{abstract}
\vfill
\noindent\small Source baseline: ProveIt revision
\texttt{\commit}.\\
The results are proposed additions to the inspected material; historical
priority beyond the cited comparison has not been established.
\end{titlepage}

\hypersetup{pageanchor=true}
\pagenumbering{roman}
\tableofcontents
\newpage
\pagenumbering{arabic}

\section{Research targets, scope, and notation}
\label{sec:scope}

The immediate starting point is the quartic digamma calculation in the
incoming reports \cite{Quartic,Barnes}.  The first higher-power question
in \cite{Quartic} asks for the fifth moment with \emph{every} finite
correction, including the correct linear Laurent coefficients and
ordinary zeta derivatives.  Its next question asks whether scalar
polynomials in the digamma function can cancel lower-depth data at
arbitrary degree.  The present article resolves the first question and
proves an all-power reduction mechanism.  It also identifies the precise
extra coordinate that must be addressed before the stronger scalar
cancellation proposal can be asserted at degree five.

The second line of work begins with the relative ordered harmonic
interpolation in the incoming independent-order report \cite{Independent}.
We prove a two-marker generating formula and exact depth cancellations
that keep a tagged complex exponent, together with all its spectral
derivatives.  This provides a useful connection between the harmonic
and polylogarithmic parts of the research program.

\begin{center}
\small
\begin{tabular}{>{\raggedright\arraybackslash}p{0.24\textwidth}
>{\raggedright\arraybackslash}p{0.39\textwidth}
>{\raggedright\arraybackslash}p{0.27\textwidth}}
\toprule
Source target & Result established here & Scope retained\\
\midrule
Fifth digamma finite part & Complete formula in Theorem~\ref{thm:q5},
with all scalar and spectral corrections & No reduction of the retained
spectral constants to a smaller arithmetic algebra\\[3pt]
Arbitrary integer powers & Normalized expansion, universal contour
constant, finite word algorithm, and ordinary subtracted sum & Fixed
integer power; no assertion for arbitrary complex powers\\[3pt]
Canonical scalar polynomials & Explicit fifth-degree cancellation of
height-one lower data; mixed derivative term remains & A formal
coordinate statement, not arithmetic independence\\[3pt]
Independent ordered exponents & Two-marker Gamma formula, integer-gap
reductions, Stieltjes and Gamma consequences & Stated convergence discs;
several tags are proved here on finite intervals\\[3pt]
Gaussian $S_6$, revised $S_8$ & A bounded expanded $S_6$ relation search
is recorded separately & Both remain conjectural\\
\bottomrule
\end{tabular}
\end{center}

The source review is targeted.  It is not a proof audit of every chapter
or every incoming archive.  The familiar Gamma recurrences and
asymptotics, finite stuffle algebra, and low-weight multiple-zeta
evaluations are antecedents, not new results \cite{DLMF,IKZ}.  Nor is the
existence of entire harmonic interpolation new: we use the established
interpolant of Ihara--Nakamura--Yamamoto \cite{INY}.  The new contribution
relative to the inspected repository is the collection of exact
extensions, normalizations, and reductions proved below.  No result is
accepted on the basis of integer-relation recognition.

\subsection{Conventions}

Let $\psi(z)=\Gamma'(z)/\Gamma(z)$, let $\gamma$ be Euler's constant,
and put
\[
 H_n^{(r)}=\sum_{j=1}^n j^{-r},\qquad H_n=H_n^{(1)},\qquad
 H_n(w)=\sum_{n\ge j_1>\cdots>j_d\ge1}
                     j_1^{-w_1}\cdots j_d^{-w_d}.
\]
The empty word has value $H_n(\varnothing)=1$; a nonempty word has
$H_0(w)=0$.  All finite words used in the digamma calculation have
positive integer entries.  We use the strict multiple-zeta convention
\begin{equation}
 Z(s_1,\ldots,s_d)=\zeta(s_1,\ldots,s_d)
 =\sum_{n_1>\cdots>n_d\ge1}
                 n_1^{-s_1}\cdots n_d^{-s_d}.
 \label{eq:strict-zeta}
\end{equation}
In particular $Z(s,w)=\sum_{n\ge1}H_{n-1}(w)n^{-s}$ where the series
converges.  At the positive integer words below, a first entry at least
two makes the series absolutely convergent.

Our Stieltjes convention is
\begin{equation}
 \zeta(1+u,a)=\frac1u+
       \sum_{j\ge0}\frac{(-1)^j}{j!}\gamma_j(a)u^j,
 \qquad a>0,\qquad \gamma_j=\gamma_j(1).
 \label{eq:stieltjes-convention}
\end{equation}
Thus $\gamma_0(a)=-\psi(a)$ and $\gamma_0=\gamma$.
For a meromorphic germ $f(u)$, $[u^j]f$ denotes its Laurent coefficient.
In particular $[u]f$ is meaningful even at a pole and is not then
$f'(0)$.  This distinction is essential in the fifth-power formula.

\begin{definition}[Original-coordinate finite part]
Suppose $f(x)=\sum_{j=1}^k a_jx^{-j}+g(x)$ near zero, with $g$
continuous on $[0,1]$.  Define
\begin{equation}
 \FP\int_0^1 f(x)\,\dd x
 =\lim_{\epsilon\downarrow0}\left\{
 \int_\epsilon^1 f(x)\,\dd x
 -\sum_{j=2}^k\frac{a_j\epsilon^{1-j}}{j-1}
 +a_1\log\epsilon\right\}.
 \label{eq:fp-definition}
\end{equation}
The coordinate is exactly $x$ and the upper endpoint is exactly $1$.
We write $Q_k=\FP\int_0^1\psi(x)^k\,\dd x$ for $k\ge1$.
\end{definition}

For a sequence $B_N=P(\log N)+o(1)$, with $P$ a polynomial, define
$\CT_{\log N}B_N=P(0)$.  This cutoff operation and a Laurent
coefficient in an auxiliary variable are not identified by convention;
Lemma~\ref{lem:cutoff} proves the comparison needed here.  Complex
logarithms in the global primitive use the principal branch on
$\C\setminus(-\infty,0]$.  All powers of positive real summation
variables use the real logarithm.

\section{Pole coefficients and the normalized global expansion}
\label{sec:poles}

The standard Taylor expansion of $\psi(1+\epsilon)$ and the recurrence
$\psi(z+1)=\psi(z)+1/z$ give \cite[\S\S5.5,5.7]{DLMF}
\begin{align}
 \psi(-n+\epsilon)
 &=-\frac1\epsilon+b_n+
                 \sum_{r\ge2}c_{r,n}\epsilon^{r-1},
 \label{eq:pole-series}\\
 b_n&=H_n-\gamma,\qquad
 c_{r,n}=H_n^{(r)}+(-1)^r\zeta(r),\qquad n\ge0.
 \notag
\end{align}
This is a convergent local expansion for $0<|\epsilon|<1$.

For each fixed positive integer $k$, define the finite polynomial data
\begin{align}
 p_{k,j}(n)
 &=[\epsilon^{k-j}]
   \left(-1+b_n\epsilon+\sum_{r=2}^k c_{r,n}\epsilon^r\right)^k,
      &&1\le j\le k, \label{eq:pdef}\\
 r_k(n)&=p_{k,1}(n),\qquad
 P_{k,n}(z)=\sum_{j=1}^k\frac{p_{k,j}(n)}{(z+n)^j},
 \qquad
 \kappa_k=[\epsilon^0]\psi(\epsilon)^k.
 \label{eq:Pdef}
\end{align}
The symbol $\kappa_k$ is a local Laurent constant; it will not be
confused with the universal contour constant $C_k$ introduced later.
The pole coefficients satisfy
\begin{equation}
 p_{k,j}(n)=O_k\bigl((1+\log(n+1))^{k-j}\bigr).
 \label{eq:p-growth}
\end{equation}
Indeed $b_n=O(1+\log(n+1))$ and every $c_{r,n}$ with $r\ge2$ is
bounded.  In a term of degree $k-j$ in $\epsilon$, at most $k-j$
factors of $b_n$ can occur.

\begin{theorem}[Normalized Mittag--Leffler formula for every power]
\label{thm:ML}
For $z\notin\{0,-1,-2,\ldots\}$,
\begin{equation}
 \boxed{\;
 \psi(z)^k=P_{k,0}(z)+\kappa_k+
       \sum_{n\ge1}\bigl(P_{k,n}(z)-P_{k,n}(0)\bigr).
 \;}
 \label{eq:ML}
\end{equation}
The series converges normally on compact sets avoiding the poles.
In particular the entire remainder is exactly the displayed constant,
with no uncomputed entire function.
\end{theorem}

\begin{proof}
Normal convergence of the candidate series follows from
\eqref{eq:p-growth}: on a fixed compact set, for sufficiently large $n$,
\[
 (z+n)^{-j}-n^{-j}=O(n^{-j-1}),
\]
uniformly in $z$.  It remains to determine the possible entire
remainder.  Put $f(w)=\psi(w)^k-P_{k,0}(w)$, so $f$ is analytic at
zero and $f(0)=\kappa_k$.  Integrate
\[
                  \frac{zf(w)}{w(w-z)}
\]
around $|w|=R_N=N+\tfrac12$, with positive orientation and $N$ large.

We need a uniform estimate on the full circle, including near the
negative real axis.  On its right half, the usual sectorial digamma
asymptotic gives $\psi(w)=O(1+\log R_N)$.  On its left half, use
\begin{equation}
             \psi(w)=\psi(1-w)-\pi\cot(\pi w).
 \label{eq:reflection}
\end{equation}
Here $1-w$ lies in the right half-plane and has size comparable to
$R_N$.  If $|\Im w|\ge1$, the cotangent is uniformly bounded.  If
$|\Im w|<1$, the points of this circle are a fixed positive distance
from the integers for all sufficiently large $N$, because its real
intersections are half-integers.  The cotangent is again uniformly
bounded.  Thus $f(w)=O_k((1+\log R_N)^k)$ on the whole contour, and
the integral is $O_k((1+\log R_N)^k/R_N)=o(1)$.

The residues at $w=z$ and $w=0$ are $f(z)$ and $-f(0)$.
At $w=-n$, the residue is
$-P_{k,n}(z)+P_{k,n}(0)$, as follows by expanding
$z/[w(w-z)]=1/(w-z)-1/w$ there.  The residue theorem and passage to
the limit give \eqref{eq:ML}.  This proves the normalization, not just
agreement of the principal parts.
\end{proof}

\subsection{A finite recurrence for the coefficient differences}

Write $\Delta p_{k,j}(n)=p_{k,j}(n)-p_{k,j}(n-1)$ for $n\ge1$,
and $\Delta r_k(n)=\Delta p_{k,1}(n)$.

\begin{proposition}[Exact difference recurrence]
\label{prop:difference}
For $1\le j\le k$,
\begin{equation}
 \Delta p_{k,j}(n)=
 \sum_{\ell=1}^{k-j}\binom{k}{\ell}
 \sum_{h=j}^{k-\ell}
       \binom{\ell+h-j-1}{h-j}
       \frac{p_{k-\ell,h}(n-1)}{n^{\ell+h-j}}.
 \label{eq:difference}
\end{equation}
An empty sum is zero; in particular $\Delta p_{k,k}(n)=0$.
\end{proposition}

\begin{proof}
The recurrence for $\psi$ gives the local identity
\[
 \psi(-n+\epsilon)=\psi(-(n-1)+\epsilon)+\frac1{n-\epsilon}.
\]
Expand its $k$th power by the binomial theorem.  In the term with
$\ell$ copies of $(n-\epsilon)^{-1}$, multiply the pole
$p_{k-\ell,h}(n-1)\epsilon^{-h}$ by
\[
 (n-\epsilon)^{-\ell}
   =\sum_{q\ge0}\binom{\ell+q-1}{q}
                       \frac{\epsilon^q}{n^{\ell+q}}.
\]
The coefficient of $\epsilon^{-j}$ has $q=h-j$, which gives
\eqref{eq:difference}.  Only finitely many coefficients contribute.
\end{proof}

The recurrence proves, in particular,
$\Delta p_{k,j}(n)=O_k((1+\log n)^k/n)$ for $j<k$.
Sharper logarithmic powers are immediate, but are not needed for the
exact identities.  Formula~\eqref{eq:difference} is also a convenient
implementation specification: it involves only lower powers and exact
rational binomial coefficients.

\section{The universal contour polynomial}
\label{sec:contour}

The subtraction constants in \eqref{eq:ML} are individually awkward
harmonic expressions.  Their combined cutoff asymptotic is much
simpler than their separate expansions suggest.

\begin{theorem}[All coefficients of the large-circle boundary term]
\label{thm:contour}
For every fixed integer $k\ge1$,
\begin{align}
 \kappa_k-\sum_{n=1}^N P_{k,n}(0)
 &=\frac1{2\pi}\int_{-\pi}^{\pi}
                         (\log N+i\theta)^k\,\dd\theta+o(1)
 \notag\\
 &=\sum_{m=0}^{\lfloor k/2\rfloor}
   \binom{k}{2m}\frac{(-1)^m\pi^{2m}}{2m+1}
                      (\log N)^{k-2m}+o(1).
 \label{eq:contour-polynomial}
\end{align}
Consequently
\begin{equation}
 C_k:=\CT_{\log N}\left(\kappa_k-\sum_{n=1}^N P_{k,n}(0)\right)
 =\begin{cases}
  (-1)^m\pi^{2m}/(2m+1),&k=2m,\\
  0,&k\text{ odd}.
 \end{cases}
 \label{eq:Ck}
\end{equation}
\end{theorem}

\begin{proof}
The residue at zero of $\psi(w)^k/w$ is $\kappa_k$.  At $-n$ it is
$-P_{k,n}(0)$.  Thus, on $R=N+\tfrac12$, the residue theorem gives
the exact identity
\begin{equation}
 \kappa_k-\sum_{n=1}^N P_{k,n}(0)
       =\frac1{2\pi}\int_{-\pi}^{\pi}
                           \psi(Re^{i\theta})^k\,\dd\theta.
 \label{eq:exact-circle}
\end{equation}
We compare this with the angular average of the principal logarithm.
On the right half-circle, $\psi(w)-\Log w=O(1/R)$ uniformly.
On the left half-circle with $|\Im w|\ge\log R$, reflection
\eqref{eq:reflection} and the exponential expansion of cotangent give
\[
 \psi(w)-\Log w=O(1/R)+O(e^{-2\pi|\Im w|}).
\]
For example, in the upper left quadrant,
$\Log(1-w)=\Log w-i\pi+O(1/R)$ and
$\pi\cot\pi w=-i\pi+O(e^{-2\pi\Im w})$; the lower quadrant is
the conjugate branch calculation.

The remaining arc near the negative real axis has angular length
$O(\log R/R)$.  On that arc, the bounded-cotangent argument in
Theorem~\ref{thm:ML} gives $\psi(w)-\Log w=O(1)$, while both
terms are $O(1+\log R)$.  Factoring the difference of the $k$th
powers shows that its contribution to the angular average is
$O_k((1+\log R)^k/R)$.  The rest of the circle contributes a quantity
of the same or smaller order.  The logarithm has a jump at the one
endpoint of the angular parametrization, which does not affect the
integral.  Therefore
\[
 \frac1{2\pi}\int_{-\pi}^{\pi}
   \bigl(\psi(Re^{i\theta})^k-(\log R+i\theta)^k\bigr)
                                      \,\dd\theta=o(1).
\]
Replacing $\log R$ by $\log N$ changes this by $o(1)$.  Odd powers
of $\theta$ integrate to zero and the even angular moments are
$(-1)^m\pi^{2m}/(2m+1)$, proving both assertions.
\end{proof}

For orientation, $C_1=C_3=C_5=0$, $C_2=-2\zeta(2)$, and
$C_4=18\zeta(4)$.  These are combined boundary constants.  They are
not the local Laurent constants $\kappa_k$.

\section{A finite-part formula for every positive integer power}
\label{sec:bridge}

Define
\begin{align}
 J_k&=\sum_{n\ge1}\sum_{j=2}^k
          \frac{\Delta p_{k,j}(n)}{(j-1)n^{j-1}},
 \label{eq:Jdef}\\
 D_k(s)&=\sum_{n\ge1}\frac{\Delta r_k(n)}{n^s},
                       \qquad \Re s>0.
 \label{eq:Ddef}
\end{align}
The first series is ordinary and absolutely convergent, by
Proposition~\ref{prop:difference}.  The second converges absolutely
in its initial half-plane.  Its meromorphic continuation to a
neighborhood of zero will follow from the proof below.

\begin{theorem}[The all-power harmonic spectral formula]
\label{thm:all-power}
For every integer $k\ge1$,
\begin{equation}
                 \boxed{\;Q_k=C_k+J_k+[s]D_k(s).\;}
 \label{eq:all-power}
\end{equation}
All data on the right are specified by the finite polynomial
calculation \eqref{eq:pdef}, the exact difference recurrence
\eqref{eq:difference}, and the coefficient of a meromorphic germ.
The identity uses the original-coordinate finite part
\eqref{eq:fp-definition}.
\end{theorem}

We first record the exact cutoff comparison that determines the sign
and the order of the spectral coefficient.

\begin{lemma}[Logarithmic cutoffs and Laurent coefficients]
\label{lem:cutoff}
Suppose
$a_n=P(\log n)+O(n^{-\delta}(1+\log n)^M)$, where $P$ is a
polynomial and $\delta>0$.  Let $A(s)=\sum_{n\ge1}a_nn^{-s}$
in its right half-plane of absolute convergence.  Then $A(1+u)$
continues meromorphically near $u=0$, and
\begin{align}
 \CT_{\log N}\sum_{n=1}^N\frac{a_n}{n}
           &=[u^0]A(1+u), \label{eq:cutoff-zero}\\
 \CT_{\log N}\sum_{n=1}^N\frac{a_n\log n}{n}
           &=-[u]A(1+u). \label{eq:cutoff-one}
\end{align}
\end{lemma}

\begin{proof}
For $P(L)=\sum_j c_jL^j$, write
\[
 A(1+u)=\sum_j c_j(-1)^j\zeta^{(j)}(1+u)
           +\sum_{n\ge1}\frac{a_n-P(\log n)}{n^{1+u}}.
\]
The last series and all its derivatives converge normally for
$\Re u>-\delta$.  The elementary Stieltjes limit formula gives
\[
 \sum_{n=1}^N\frac{(\log n)^j}{n}
       =\frac{(\log N)^{j+1}}{j+1}+\gamma_j
                    +O((1+\log N)^j/N).
\]
The stated error follows by one sum--integral comparison applied
to the derivative of $(\log x)^j/x$ on the tail.
On the Laurent side,
$[u^0](-1)^j\zeta^{(j)}(1+u)=\gamma_j$ and
$[u](-1)^j\zeta^{(j)}(1+u)=-\gamma_{j+1}$.
Termwise comparison proves the first formula and, after inserting
one further factor of $\log n$, the second.  Absolute convergence
handles the subtracted remainder in both cases.
\end{proof}

\begin{proof}[Proof of Theorem~\ref{thm:all-power}]
The harmonic-number asymptotic expansions, or reflection followed by
the positive-axis digamma expansion, show that
\begin{align}
 r_k(n)&=R_k(\log n)+O_k((1+\log n)^M/n),\label{eq:r-asymptotic}\\
 \Delta r_k(n)&=\frac{R_k'(\log n)}n
                  +O_k((1+\log n)^M/n^2),\label{eq:delta-asymptotic}\\
 R_k(L)&=\Res_{\epsilon=0}(L-\pi\cot\pi\epsilon)^k.
 \label{eq:Rdefinition}
\end{align}
Here and below $M$ may depend on $k$.  To justify the second estimate,
use the complete harmonic asymptotic expansions through the $1/n$
term: $r_k(n)=R_k(\log n)+n^{-1}U_k(\log n)
+O((1+\log n)^M/n^2)$ with a polynomial $U_k$.  The difference of
the explicit $n^{-1}U_k$ terms is $O((1+\log n)^M/n^2)$, and
the leading difference is $R_k'(\log n)/n$ to that same accuracy.
Thus the second formula does not follow from differencing an
unspecified first-order error.  It also follows directly from
\eqref{eq:difference} and the first two harmonic asymptotic terms.

Apply Lemma~\ref{lem:cutoff} to $a_n=n\Delta r_k(n)$.
It proves the meromorphic continuation of $D_k$ and gives
\begin{equation}
 \CT_{\log N}\sum_{n=1}^N\Delta r_k(n)\log n=-[s]D_k(s).
 \label{eq:delta-cutoff}
\end{equation}

Now integrate the normally convergent series \eqref{eq:ML}.
After the pole $P_{k,0}$ is removed, the same
$O(n^{-j-1})$ estimate proves uniform convergence also on the
closed interval $[0,1]$, so integration across its regular
endpoint is justified.
The finite part of the pole at zero contributes
$-\sum_{j=2}^kp_{k,j}(0)/(j-1)$; its simple pole has finite part
zero.  For a pole at $-n$, the integrated simple-pole term is
$r_k(n)\log(1+1/n)$, and the higher-pole terms have numerators
$p_{k,j}(n)(n^{1-j}-(n+1)^{1-j})/(j-1)$.
The constant subtraction integrates to $-P_{k,n}(0)$.

Use the following \emph{finite} summation identities:
\begin{align}
 &\sum_{n=1}^N p_{k,j}(n)
               \bigl(n^{1-j}-(n+1)^{1-j}\bigr)\notag\\
 &\hspace{6mm}=p_{k,j}(0)+
           \sum_{n=1}^N\frac{\Delta p_{k,j}(n)}{n^{j-1}}
                    -\frac{p_{k,j}(N)}{(N+1)^{j-1}},
                  &&j\ge2,\label{eq:sbp-high}\\
 &\sum_{n=1}^N r_k(n)\log(1+1/n)
     =r_k(N)\log(N+1)-\sum_{n=1}^N\Delta r_k(n)\log n.
 \label{eq:sbp-simple}
\end{align}
The first terms in \eqref{eq:sbp-high} cancel the finite part of
the pole at zero exactly, while its final terms tend to zero.
If $J_{k,N}$ denotes the cutoff in \eqref{eq:Jdef}, we obtain
\begin{align}
 Q_k=\lim_{N\to\infty}\biggl\{
 &\kappa_k-\sum_{n=1}^N P_{k,n}(0)+J_{k,N}
       +r_k(N)\log(N+1)\notag\\
 &\hspace{12mm}-\sum_{n=1}^N\Delta r_k(n)\log n\biggr\}.
 \label{eq:Qcutoff}
\end{align}
All terms here arose from finite sums or normally convergent series.
By \eqref{eq:r-asymptotic}, the polynomial asymptotic of the boundary
$r_k(N)\log(N+1)$ is $(\log N)R_k(\log N)$, which has no constant
term.  Theorem~\ref{thm:contour} supplies the constant $C_k$ for the
first two terms of \eqref{eq:Qcutoff}, and \eqref{eq:delta-cutoff}
supplies $+[s]D_k$.  Finally $J_{k,N}\to J_k$.  Taking the constant
coefficient of the complete cutoff expression proves
\eqref{eq:all-power}.
\end{proof}

\subsection{A closed ordinary-series form}

The residue polynomial in \eqref{eq:Rdefinition} itself has a closed
form.  This makes every Stieltjes correction explicit, without naming
any multiple-zeta Laurent coefficient.

\begin{lemma}[Closed asymptotic residue polynomial]
\label{powers:lem:closed-residue}
For an integer \(k\ge0\), define
\[
 R_k(L)=\Res_{t=0}\bigl(L-\pi\cot(\pi t)\bigr)^k.
\]
Then
\begin{equation}
 \boxed{\displaystyle
 R_k(L)=-\frac{(L+i\pi)^k-(L-i\pi)^k}{2i\pi}.}
 \label{powers:eq:closed-residue}
\end{equation}
For \(k\ge1\), if \(R_k'(L)=\sum_{j\ge0}a_{k,j}L^j\), the
ordinary Stieltjes contribution is the explicit finite expression
\begin{equation}
 \boxed{\displaystyle
 -\sum_{j\ge0}a_{k,j}\gamma_{j+1}
 =k\sum_{m=0}^{\lfloor(k-2)/2\rfloor}
       (-1)^m\binom{k-1}{2m+1}\pi^{2m}\gamma_{k-2m-1}.}
 \label{powers:eq:closed-Stieltjes}
\end{equation}
The right side is understood as zero for \(k=1\). Here
\(\zeta(1+s)=s^{-1}+\sum_{j\ge0}(-1)^j\gamma_j s^j/j!\).
\end{lemma}

\begin{proof}
Put \(f(t)=L-\pi\cot(\pi t)\). The elementary differential
identity
\[
 f'(t)=(f(t)-L)^2+\pi^2
\]
and the zero residue of a meromorphic derivative give, for \(k\ge2\),
\[
 0=\frac1{k-1}\Res_{t=0}\frac{\dd}{\dd t}f(t)^{k-1}
   =R_k(L)-2LR_{k-1}(L)+(L^2+\pi^2)R_{k-2}(L).
\]
Since \(R_0=0\) and \(R_1=-1\), solving this recurrence with
characteristic roots \(L\pm i\pi\) proves
\eqref{powers:eq:closed-residue}. Expanding its numerator yields
\[
 R_k(L)=-\sum_{m=0}^{\lfloor(k-1)/2\rfloor}
       (-1)^m\binom{k}{2m+1}\pi^{2m}L^{k-2m-1}.
\]
On differentiating in \(L\), the possible term with
\(k-2m-1=0\) vanishes. Use
\((k-2m-1)\binom{k}{2m+1}=k\binom{k-1}{2m+1}\).
This gives \eqref{powers:eq:closed-Stieltjes}.
\end{proof}

To see directly why this finite expression is the Stieltjes term,
suppose the residue difference satisfies the previously established
asymptotic formula
\[
 r_k(n)-r_k(n-1)
   =\frac{R_k'(\log n)}n
          +O\!\left(\frac{(1+\log n)^M}{n^2}\right).
\]
For \(D_k(s)=\sum_{n\ge1}(r_k(n)-r_k(n-1))n^{-s}\), subtraction
therefore gives
\[
 D_k(s)=\sum_{j\ge0}a_{k,j}(-1)^j\zeta^{(j)}(1+s)
 +\sum_{n\ge1}
   \left(r_k(n)-r_k(n-1)-\frac{R_k'(\log n)}n\right)n^{-s}.
\]
The second series and its first spectral derivative converge normally
near \(s=0\). Since
\([s](-1)^j\zeta^{(j)}(1+s)=-\gamma_{j+1}\), it follows that
\begin{align}
 [s]D_k(s)
 ={}&k\sum_{m=0}^{\lfloor(k-2)/2\rfloor}
       (-1)^m\binom{k-1}{2m+1}\pi^{2m}\gamma_{k-2m-1}
       \notag\\
 &-\sum_{n\ge1}
   \left(r_k(n)-r_k(n-1)-\frac{R_k'(\log n)}n\right)\log n.
 \label{powers:eq:closed-Dlinear}
\end{align}
The final sum is absolutely convergent. For example, the finite
Stieltjes terms at powers four and five are respectively
\(12\gamma_3-24\zeta(2)\gamma_1\) and
\(20\gamma_4-120\zeta(2)\gamma_2\).


\begin{corollary}[Ordinary subtracted series at every power]
\label{cor:ordinary-all}
With the definitions above,
\begin{align}
 Q_k={}&C_k+J_k+
 k\sum_{m=0}^{\lfloor(k-2)/2\rfloor}
       (-1)^m\binom{k-1}{2m+1}\pi^{2m}\gamma_{k-2m-1}
 \notag\\
 &-\sum_{n\ge1}
   \left(\Delta r_k(n)-\frac{R_k'(\log n)}n\right)\log n.
 \label{eq:ordinary-all}
\end{align}
The finite Stieltjes sum is zero at $k=1$.  Both infinite sums are
ordinary absolutely convergent sums.  The last summand is
$O_k((1+\log n)^{M+1}/n^2)$.
\end{corollary}

This is useful both conceptually and computationally.  The right side
contains no unspecified regularization and no uncomputed entire
remainder.  Further terms of the same harmonic asymptotic expansions
give convergent Hurwitz-zeta tail accelerations if numerical evaluation
is wanted.  Such acceleration is optional for the identity itself.

\section{A finite algorithm in strict harmonic words}
\label{sec:algorithm}

The spectral formula becomes finite algebra once its target class is
specified.  Work over $\Q[\gamma,\zeta(2),\ldots,\zeta(k)]$ and put
$h_r=H_{n-1}^{(r)}$, $v=1/n$.  Compute the polynomial
$p_{k,j}$ from \eqref{eq:pdef}, and replace $H_n^{(r)}$ by
$h_r+v^r$ to obtain its value at $n$.  Subtract its value at $n-1$.
This produces every $\Delta p_{k,j}$ as a finite polynomial in
$v,h_1,\ldots,h_k$.

The finite stuffle recursion, with $(a,A)$ denoting a leading letter
followed by a word, is
\begin{equation}
 (a,A)*(b,B)=(a,A*(b,B))+(b,(a,A)*B)+(a+b,A*B).
 \label{eq:stuffle}
\end{equation}
It expresses products of the $h_r$ as a finite integer combination
of strict words $H_{n-1}(w)$.  The coefficients count equal-index
possibilities as well as the two strict orderings.  Thus
\[
        h_1^2=2H_{n-1}(1,1)+H_{n-1}^{(2)},
\]
and a further multiplication gives
\[
 h_1^3=6H_{n-1}(1,1,1)+3H_{n-1}(2,1)
                  +3H_{n-1}(1,2)+H_{n-1}^{(3)}.
\]
These are finite identities before any infinite sum is taken.

Every positive-integer word has a polynomial harmonic asymptotic
\begin{equation}
 H_{n-1}(w)=B_w(\log n)+O((1+\log n)^M/n)
 \label{eq:word-asymptotic}
\end{equation}
for some polynomial $B_w$ and integer $M$.  Here is a direct induction
that also supplies the continuation needed term by term.  For the
empty word the assertion is immediate.  Given it for $w$, use
$H_N(a,w)=\sum_{m=1}^N H_{m-1}(w)m^{-a}$.  If $a=1$,
the Stieltjes cutoff formula applied to each power of $\log m$
and the absolutely convergent subtracted remainder give a
polynomial in $\log N$ with an
$O((1+\log N)^{M'}/N)$ error.  If $a\ge2$, the sum converges
and its tail has that same error form, so the polynomial is a
constant.  Replacing $N$ by $n-1$ preserves the claimed form.
Consequently each $Z(1+s,w)$ has a meromorphic germ at zero by
Lemma~\ref{lem:cutoff}; for first entry at least two its defining
series is already holomorphic near zero.

\begin{theorem}[Terminating harmonic-word reduction]
\label{thm:algorithm}
At a fixed positive integer power $k$, the following procedure terminates
and yields the exact formula \eqref{eq:all-power}:
\begin{enumerate}
\item Compute $p_{k,j}$ and their finite differences by
\eqref{eq:pdef} or \eqref{eq:difference}.
\item Expand the summand of $J_k$ and the polynomial $\Delta r_k(n)$
in $v,h_1,\ldots,h_k$.
\item Apply the finite stuffle recursion to every harmonic monomial.
\item Replace $v^qH_{n-1}(w)$ by $Z(q,w)$ for $J_k$, and by
$Z(q+s,w)$ for $D_k(s)$.
\end{enumerate}
Every first entry in $J_k$ is at least two.  Every first entry in
$D_k(0)$ is at least one.  The resulting word weights are at most
$k-1$; coefficient factors in $\gamma$ and ordinary zeta values account
for the remaining formal weight.  The prescription introduces only a
finite number of specified linear Laurent coefficients.
\end{theorem}

\begin{proof}
Coefficient extraction and finite differences are finite polynomial
operations.  Each application of the stuffle recursion reduces the
combined remaining word length, so it terminates.  By
\eqref{eq:difference}, every monomial in $\Delta p_{k,j}$ contains
at least one $v$, and division by $n^{j-1}$ gives $q\ge j\ge2$
in $J_k$.  For $\Delta r_k$, it gives $q\ge1$.  Therefore all
replacements in the initial half-planes are identities of absolutely
convergent series.  Their continuations at zero agree by the
meromorphic subtraction justified by \eqref{eq:word-asymptotic}.
Finally assign weight $1$ to $h_1,\gamma,v$, weight $r$ to
$h_r,\zeta(r)$, and use \eqref{eq:pdef}.  Both relevant summand
polynomials have weight $k-1$, and stuffle preserves it.
\end{proof}

This is a reduction to a precise finite class of multiple-zeta values
and their spectral coefficients.  It does not prove that all those
values reduce further to products of ordinary zeta values.  For higher
powers it is appropriate to retain admissible multiple-zeta words
unless an exact additional relation is supplied.

\subsection{Recovery of the lower-power formulas}

Define the height-one coefficients
\begin{equation}
 \eta=[u]Z(1+u,1),\qquad
 \theta=[u]Z(1+u,1,1),\qquad
 \chi=[u]Z(1+u,1,1,1).
 \label{eq:height-coefficients}
\end{equation}
The first instances of Theorem~\ref{thm:all-power} are
\begin{align}
 Q_1={}&0,\notag\\
 Q_2={}&2\gamma_1-2\zeta(2),\label{eq:Q2}\\
 Q_3={}&3\zeta(2)-6\eta-6\gamma\gamma_1,\label{eq:Q3}\\
 Q_4={}&12\zeta(3)-12\gamma\zeta(2)+18\zeta(4)+4\zeta'(3)
       -24\theta+24\gamma\eta
                  +12(\gamma^2-\zeta(2))\gamma_1.
 \label{eq:Q4}
\end{align}
The cubic and quartic identities are consistency checks against the
incoming reports \cite{Quartic,Barnes}; they are not presented as new
evaluations here.  For example,
\begin{align*}
 D_2(s)&=-2\zeta(1+s),& J_2&=0,\\
 D_3(s)&=-6Z(1+s,1)+6\gamma\zeta(1+s),&J_3&=3\zeta(2),\\
 D_4(s)&=-24Z(1+s,1,1)+24\gamma Z(1+s,1)\\
       &\quad+12(\zeta(2)-\gamma^2)\zeta(1+s)+4\zeta(3+s),
       &J_4&=12\zeta(3)-12\gamma\zeta(2).
\end{align*}
At $k=1$ the residue is the constant $-1$, hence $J_1=D_1=0$;
the finite part also follows directly from $\log\Gamma(1)=0$ and
$\log\Gamma(\epsilon)=-\log\epsilon+O(\epsilon)$.

\section{The fifth power with every finite correction}
\label{sec:fifth}

For compactness write $\zeta_r=\zeta(r)$ in this section.  At a fixed
pole use $b=b_n$, $c=c_{2,n}$, $d=c_{3,n}$, and $e=c_{4,n}$.
Direct coefficient extraction gives
\begin{align}
 p_{5,5}&=-1,\qquad
 p_{5,4}=5b,\qquad
 p_{5,3}=5c-10b^2,\notag\\
 p_{5,2}&=10b^3-20bc+5d,\notag\\
 r_5&=-5b^4+30b^2c-10c^2-20bd+5e.
 \label{eq:p5}
\end{align}
For reference, the local constant is
\begin{equation}
 \kappa_5=-\gamma^5+20\gamma^3\zeta_2-55\gamma\zeta_4
             -30\gamma^2\zeta_3+20\zeta_2\zeta_3-5\zeta_5.
 \label{eq:kappa5}
\end{equation}
The universal constant is $C_5=0$, a different quantity.

For the finite differences, let $b,c,d$ instead denote the preceding
index $n-1$.  Then
\begin{align}
 \Delta p_{5,4}(n)&=\frac5n,\notag\\
 \Delta p_{5,3}(n)&=-\frac5{n^2}-\frac{20b}{n},\notag\\
 \Delta p_{5,2}(n)&=-\frac5{n^3}+\frac{10b}{n^2}
                              +\frac{30b^2-20c}{n},\notag\\
 \Delta r_5(n)&=\frac{-20b^3+60bc-20d}{n}
                          +\frac{10c}{n^2}+\frac{20b}{n^3}.
 \label{eq:delta5}
\end{align}
Equivalently the first numerator in the last line is $5r_4(n-1)$.
These identities can be checked simply by replacing
$(b,c,d,e)$ by $(b+1/n,c+1/n^2,d+1/n^3,e+1/n^4)$ in
\eqref{eq:p5}; Proposition~\ref{prop:difference} gives an independent
binomial derivation.

The higher poles therefore contribute
\begin{align}
 J_5
 &=\sum_{n\ge1}\left\{
       \frac{30(H_{n-1}-\gamma)^2
                    -20(H_{n-1}^{(2)}+\zeta_2)}{n^2}
                            -\frac{35}{6n^4}\right\}\notag\\
 &=30\gamma^2\zeta_2-60\gamma\zeta_3+\frac{35}{3}\zeta_4.
 \label{eq:J5}
\end{align}
For completeness, the low-weight evaluations needed in this last step
are
\[
 Z(2,1)=\zeta_3,\quad Z(2,1,1)=\zeta_4,\quad
 Z(2,2)=\frac34\zeta_4,\quad \zeta_2^2=\frac52\zeta_4.
\]
They follow from the standard duality, sum, and stuffle identities
\cite{DLMF,IKZ}.  In particular,
\begin{align*}
 \sum_{n\ge1}\frac{(H_{n-1}-\gamma)^2}{n^2}
       &=\frac{11}{4}\zeta_4-2\gamma\zeta_3+\gamma^2\zeta_2,\\
 \sum_{n\ge1}\frac{H_{n-1}^{(2)}+\zeta_2}{n^2}
       &=\frac{13}{4}\zeta_4,
\end{align*}
which proves every scalar term in \eqref{eq:J5}.

Applying finite stuffle to the last line of \eqref{eq:delta5} yields
the complete spectral identity
\begin{align}
 D_5(s)={}&-120Z(1+s,1,1,1)+120\gamma Z(1+s,1,1)
 \notag\\
 &+60(\zeta_2-\gamma^2)Z(1+s,1)
       +20(\gamma^3-3\gamma\zeta_2+\zeta_3)\zeta(1+s)
 \notag\\
 &+20Z(1+s,3)+10Z(2+s,2)+20Z(3+s,1)
 \notag\\
 &+10\zeta_2\zeta(2+s)-20\gamma\zeta(3+s).
 \label{eq:D5}
\end{align}
All equalities here hold first for $\Re s>0$, and then as meromorphic
germs.  In the first numerator of \eqref{eq:delta5}, the words
$(2,1)$ and $(1,2)$ from $H_{n-1}^3$ cancel against the corresponding
words in $H_{n-1}H_{n-1}^{(2)}$.  The residual $(3)$ word does not
cancel.  It is responsible for $Z(1+s,3)$.

Set
\begin{equation}
 D_{p,q}=[u]Z(p+u,q),\qquad
 \Omega_4=2D_{1,3}+D_{2,2}+2D_{3,1}.
 \label{eq:mixed-def}
\end{equation}
For $p\ge2$ these coefficients are ordinary outer-order derivatives,
$D_{p,q}=\partial_1Z(p,q)$.  For $p=1,q=3$, coefficient extraction
is required at a simple pole.

\begin{theorem}[Complete fifth-power finite part]
\label{thm:q5}
In the conventions \eqref{eq:fp-definition},
\eqref{eq:stieltjes-convention}, \eqref{eq:height-coefficients}, and
\eqref{eq:mixed-def},
\begin{align}
 Q_5={}&30\gamma^2\zeta_2-60\gamma\zeta_3+\frac{35}{3}\zeta_4
          -120\chi+120\gamma\theta
          +60(\zeta_2-\gamma^2)\eta
 \notag\\
 &-20(\gamma^3-3\gamma\zeta_2+\zeta_3)\gamma_1
          +20D_{1,3}+10D_{2,2}+20D_{3,1}
 \notag\\
 &+10\zeta_2\zeta'(2)-20\gamma\zeta'(3).
 \label{eq:Q5}
\end{align}
Equivalently the three mixed terms are $10\Omega_4$.
\end{theorem}

\begin{proof}
Use $C_5=0$, \eqref{eq:J5}, and the linear coefficient of
\eqref{eq:D5} in Theorem~\ref{thm:all-power}.  In particular
$[s]\zeta(1+s)=-\gamma_1$, which fixes the sign of its correction.
The other coefficients are exactly those named in
\eqref{eq:height-coefficients} and \eqref{eq:mixed-def}.
Every step rests on a finite algebraic identity or the normalized
analytic results already proved.
\end{proof}

This answers the fifth-power research question in \cite{Quartic}:
the fourfold height-one coefficient enters with multiplier $-120$,
and the mixed twofold coefficients and both ordinary zeta derivatives
are retained explicitly.  The result is not a reduction of $\chi$
or $\Omega_4$ to ordinary zeta and Stieltjes constants.

\subsection{An ordinary convergent formula for the same finite part}

The residue polynomial and its derivative are
\[
 R_5(L)=-5L^4+60\zeta_2L^2-90\zeta_4,
 \qquad R_5'(L)=-20L^3+120\zeta_2L.
\]
Consequently Corollary~\ref{cor:ordinary-all} gives the exact alternative
\begin{align}
 Q_5={}&30\gamma^2\zeta_2-60\gamma\zeta_3+\frac{35}{3}\zeta_4
                   +20\gamma_4-120\zeta_2\gamma_2
 \notag\\
 &-\sum_{n\ge1}
  \left\{\Delta r_5(n)
           +\frac{20\log^3n-120\zeta_2\log n}{n}\right\}\log n.
 \label{eq:Q5ordinary}
\end{align}
Here $\Delta r_5$ is the explicit harmonic expression
\eqref{eq:delta5}, with $b,c,d$ at $n-1$.
It can equally be written in polygamma values at positive integers:
\begin{equation}
 b=\psi(n),\qquad c=2\zeta_2-\psi'(n),\qquad
 d=\tfrac12\psi''(n).
 \label{eq:polygamma5}
\end{equation}
The final series in \eqref{eq:Q5ordinary} converges absolutely;
its summand is $O((1+\log n)^4/n^2)$.
This gives a definition-free numerical route to $Q_5$ using only
ordinary Stieltjes constants and ordinary polygamma evaluations,
with no multivariate analytic continuation routine.

\subsection{Explicit kernels for the height-one coefficients}
\label{subsec:heightkernel}

The retained height-one coefficient has a concrete normalized
integral representation.  We give it at every depth so that the
fourfold coefficient is specified independently of the digamma moment.
Height-one Laurent constants and related normalized integrals have
antecedents in Young's work \cite[Section~4]{Young}.  The proposition
below records a direct Mellin subtraction with the coefficient and
Stieltjes conventions fixed explicitly for this article.
Let
\begin{equation}
 \frac1{\Gamma(1+u)}=\sum_{j\ge0}a_ju^j,\qquad
 K_d(t)=\frac{[-\log(1-e^{-t})]^{d-1}}
                       {(d-1)!(e^t-1)},\quad d\ge1.
 \label{eq:Kd}
\end{equation}

\begin{proposition}[Normalized height-one Mellin coefficient]
\label{prop:heightkernel}
For $d\ge1$,
\begin{equation}
 Z(1+u,1^{d-1})=\sum_{j=0}^d a_ju^{j-d}+L_du+O(u^2),
 \label{eq:heightLaurent}
\end{equation}
where
\begin{align}
 L_d=a_{d+1}
 &+\int_0^1\log t\left\{K_d(t)
                  -\frac{(-\log t)^{d-1}}{(d-1)!t}\right\}\,\dd t
 \notag\\
 &+\int_1^\infty\log t\,K_d(t)\,\dd t.
 \label{eq:heightkernel}
\end{align}
Both integrals are absolutely convergent.  In particular
$L_2=\eta$, $L_3=\theta$, and $L_4=\chi$.
\end{proposition}

\begin{proof}
The generating series of $H_{n-1}(1^{d-1})$ gives, for $\Re u>0$,
\[
 \Gamma(1+u)Z(1+u,1^{d-1})=\int_0^\infty t^uK_d(t)\,\dd t.
\]
For example it follows by dividing
$\Li_{1^{d-1}}(z)=(-\log(1-z))^{d-1}/(d-1)!$
by $1-z$, multiplying by $z$, and setting $z=e^{-t}$; for $d=1$
the kernel is simply $1/(e^t-1)$.  Subtract the leading kernel on
$(0,1)$ to write this integral as $u^{-d}+B_d(u)$, with $B_d$
holomorphic near zero.  The subtraction is integrable with any fixed
power of $\log t$, because its remainder is at worst a polynomial in
$|\log t|$ near zero; the kernel decays exponentially at infinity.

The derivative identity
\[
 K_d(t)=-\frac1{d!}\frac{\dd}{\dd t}
                         [-\log(1-e^{-t})]^d
\]
shows that $B_d(0)=0$: the integrals from $\epsilon$ to infinity
and of the subtracted kernel from $\epsilon$ to $1$ have difference
\[
 \frac{[-\log(1-e^{-\epsilon})]^d-(-\log\epsilon)^d}{d!}
                         \longrightarrow0.
\]
Multiplication by $1/\Gamma(1+u)$ now proves
\eqref{eq:heightLaurent}, with $L_d=a_{d+1}+B_d'(0)$ exactly as
in \eqref{eq:heightkernel}.
\end{proof}

For the fifth moment the coefficient used in this kernel is explicitly
\begin{equation}
 a_5=\frac{\gamma^5-10\gamma^3\zeta_2+20\gamma^2\zeta_3
       +15\gamma\zeta_2^2-30\gamma\zeta_4
       -20\zeta_2\zeta_3+24\zeta_5}{120}.
 \label{eq:a5}
\end{equation}
This follows by exponentiating
$\log(1/\Gamma(1+u))=\gamma u+
\sum_{r\ge2}(-1)^{r+1}\zeta(r)u^r/r$.
The kernel is an exact analytic representation, not by itself a
reduction of the arithmetic content of $\chi$.

\subsection{The mixed derivative and its polylogarithmic representations}
\label{subsec:mixed}

The stuffle identity in a right half-plane gives
\[
 Z(1+u,3)=\zeta(1+u)\zeta_3-Z(3,1+u)-\zeta(4+u).
\]
Consequently
\begin{equation}
 D_{1,3}=-\gamma_1\zeta_3-\partial_2Z(3,1)-\zeta'(4).
 \label{eq:D13inner}
\end{equation}
The inner-order derivative has an absolutely convergent defining
series.  In particular this transformation replaces a singular
outer-order coefficient by a convergent inner-order derivative; it
does not remove the nonsymmetric spectral information.

Define the Tornheim function
\[
 T(a,b,c)=\sum_{m,n\ge1}\frac1{m^an^b(m+n)^c}
\]
in its absolute-convergence domain.  A finite partial fraction
calculation gives, for $u$ near zero,
\begin{equation}
 T(2,2,u)=2Z(2+u,2)+4Z(3+u,1).
 \label{eq:T22}
\end{equation}
Indeed, for $l=m+n$, summation over $1\le m<l$ uses
\[
 \frac1{m^2(l-m)^2}
 =\frac1{l^2}\left(\frac1{m^2}+\frac1{(l-m)^2}
                    +\frac2{m(l-m)}\right)
\]
and gives $2H_{l-1}^{(2)}/l^2+4H_{l-1}/l^3$.
Thus
\begin{align}
 \Omega_4
 &=2D_{1,3}+\tfrac12\partial_3T(2,2,0)
 \label{eq:OmegaT}\\
 &=2(\partial_1-\partial_2)Z(3,1)+\partial_1Z(2,2)
                      -2\gamma_1\zeta_3-2\zeta'(4).
 \label{eq:OmegaInner}
\end{align}
Spectral derivatives of Tornheim sums have substantial prior literature,
including \cite{BB2014}.  These equalities are useful coordinate changes;
no arithmetic independence is asserted for the displayed coordinates.

There are also completely convergent one-dimensional representations:
\begin{align}
 D_{1,3}
  &=-\gamma_1\zeta_3+
        \sum_{n\ge1}\frac{\zeta(3,n)\log n}{n},
 \label{eq:D13tail}\\
 \partial_3T(2,2,0)
  &=\int_0^\infty
        \frac{\Li_2(e^{-t})^2-e^{-t}\zeta_2^2}{t}\,\dd t.
 \label{eq:Tpolylog}
\end{align}
For \eqref{eq:D13tail}, use
$H_{n-1}^{(3)}=\zeta_3-\zeta(3,n)$ and extract the linear
coefficient.  For \eqref{eq:Tpolylog}, insert the Frullani identity
$-\log(m+n)=\int_0^\infty(e^{-(m+n)t}-e^{-t})\,\dd t/t$
and sum.  Absolute integrability follows either from
$\sum_{m,n\ge1}\log(m+n)/(m^2n^2)<\infty$, or directly from
the $O(\log t)$ integrand at zero and exponential decay at infinity.
Equations~\eqref{eq:OmegaT}, \eqref{eq:D13tail}, and
\eqref{eq:Tpolylog} connect the first mixed correction in a digamma
power to a dilogarithm integral and a Hurwitz-zeta tail sum.

\subsection{What scalar polynomial cancellation proves}
\label{subsec:canonical}

\begin{corollary}[A normalized fifth-degree cancellation]
\label{powers:cor:canonical-fifth}
Use the fixed-coordinate finite parts
\(Q_k=\FP\int_0^1\psi(x)^k\,\dd x\), and put
\[
 \Omega_4=2D_{1,3}+D_{2,2}+2D_{3,1},
 \qquad D_{p,q}=[s]\zeta(p+s,q).
\]
Here \([s]\) denotes the coefficient of \(s\) in the Laurent
expansion, including when a pole is present. Define
\begin{align}
 P_5(T)={}&T^5+5\gamma T^4
      +10\bigl(\gamma^2+\zeta(2)\bigr)T^3\notag\\
 &\quad+10\bigl(\gamma^3+3\gamma\zeta(2)+\zeta(3)\bigr)T^2.
 \label{powers:eq:canonical-fifth-polynomial}
\end{align}
Then
\begin{align}
 \FP\int_0^1P_5\bigl(\psi(x)\bigr)\,\dd x
 ={}&-120\chi+10\Omega_4+10\zeta(2)\zeta'(2)\notag\\
 &+\frac{260}{3}\zeta(4)-60\gamma\zeta(4)
      -20\gamma^3\zeta(2)-20\zeta(2)\zeta(3).
 \label{powers:eq:canonical-fifth-finite-part}
\end{align}
Among monic polynomials
\(T^5+a_4T^4+a_3T^3+a_2T^2\) with coefficients in
\(\mathbb Q[\gamma,\zeta(2),\zeta(3)]\), the displayed polynomial is
the unique one whose coefficients of \(\theta,\eta,\gamma_1\)
vanish in the stated formulas for \(Q_2,Q_3,Q_4,Q_5\), when these
symbols are treated as formal coordinates.
\end{corollary}

\begin{proof}
The finite-part operation in one fixed coordinate is linear. In
\(Q_5+a_4Q_4+a_3Q_3+a_2Q_2\), cancellation of the displayed
\(\theta,\eta,\gamma_1\) coefficients gives, successively,
\begin{align*}
 120\gamma-24a_4&=0,\\
 60\bigl(\zeta(2)-\gamma^2\bigr)+24\gamma a_4-6a_3&=0,\\
 -20\bigl(\gamma^3-3\gamma\zeta(2)+\zeta(3)\bigr)
 +12\bigl(\gamma^2-\zeta(2)\bigr)a_4-6\gamma a_3+2a_2&=0.
\end{align*}
This triangular system has exactly the coefficients in
\eqref{powers:eq:canonical-fifth-polynomial}. The
\(\zeta'(3)\) coefficient simultaneously becomes
\(-20\gamma+4a_4=0\), while the coefficients of \(\chi\),
\(\Omega_4\), and \(\zeta(2)\zeta'(2)\) remain respectively
\(-120,10,10\). Before using any even-zeta reduction, the remaining
scalar is
\[
 \frac{35}{3}\zeta(4)+90\gamma\zeta(4)+30\zeta(2)^2
 -20\gamma^3\zeta(2)-60\gamma\zeta(2)^2
 -20\zeta(2)\zeta(3).
\]
Substitute \(\zeta(2)^2=\tfrac52\zeta(4)\) to obtain
\eqref{powers:eq:canonical-fifth-finite-part}.
\end{proof}

The uniqueness assertion is deliberately relative to the displayed
coefficient presentation and to the normalization excluding constant
and linear terms. It does not assert arithmetic independence of
\(\chi,\Omega_4,\theta,\eta,\gamma_1\), or the impossibility of
rewriting \(\Omega_4\) in another family of constants. In particular,
``cancellation of \(\gamma_1\)'' means cancellation of its explicit
coefficient while \(\Omega_4\) is retained as a coordinate; a
different representation of \(\Omega_4\) may contain
\(\gamma_1\). Further arithmetic reduction of the mixed-depth term
remains a separate question.


\section{Normalized primitives and polygamma identities}
\label{sec:primitives}

The same normalization gives a primitive, not only a finite moment.
On $\C\setminus(-\infty,0]$ use the principal logarithm.

\begin{theorem}[An all-power primitive with fixed normalization]
\label{thm:primitive}
For $k\ge1$, put
\begin{align}
 \mathcal A_k(z)={}&\kappa_k z+p_{k,1}(0)\Log z
             +\sum_{j=2}^k\frac{p_{k,j}(0)z^{1-j}}{1-j}
 \notag\\
 &+\sum_{n\ge1}\biggl\{
       p_{k,1}(n)\log(1+z/n)-zP_{k,n}(0)
 \notag\\
 &\hspace{26mm}
       +\sum_{j=2}^k p_{k,j}(n)
             \frac{(z+n)^{1-j}-n^{1-j}}{1-j}\biggr\}.
 \label{eq:primitive}
\end{align}
The series converges normally on compact subsets of that slit plane,
and
\begin{equation}
             \mathcal A_k'(z)=\psi(z)^k,\qquad
             \mathcal A_k(1)=Q_k.
 \label{eq:primitive-properties}
\end{equation}
Its local Laurent--log expansion at zero has constant coefficient
zero, which fixes the otherwise arbitrary additive constant.
\end{theorem}

\begin{proof}
Each brace in \eqref{eq:primitive} is the integral from zero to $z$
of $P_{k,n}(w)-P_{k,n}(0)$.  For large $n$ and $z$ in a fixed
compact set its magnitude is
$O_k((1+\log n)^{k-1}/n^2)$, as follows by cancellation of its
linear terms or from the bound in the proof of Theorem~\ref{thm:ML}.
The normally convergent series may therefore be differentiated, and
\eqref{eq:ML} gives the first identity.  Every term in the braces
vanishes at zero and the analytic part begins with $\kappa_kz$.
The Laurent and logarithmic terms outside the braces also have no
constant coefficient.  The definition of the original-coordinate
finite part now gives $\mathcal A_k(1)=Q_k$.
\end{proof}

For $k=1$, this is the normalized branch of $\log\Gamma(z)$,
whose local expansion begins $-\Log z-\gamma z+\cdots$.
For every higher power it gives an explicitly normalized antiderivative
through harmonic pole coefficients.  Ordinary integrals along paths
contained in the slit plane follow by endpoint subtraction of
$\mathcal A_k$.  On other pole-free intervals the same conclusion
holds after choosing a consistent analytic continuation.

\begin{corollary}[Differentiated polygamma identities]
\label{cor:derivatives}
For $m\ge1$ and $z$ away from the nonpositive integers,
\begin{equation}
 \frac{\dd^m}{\dd z^m}\psi(z)^k
  =(-1)^m\sum_{n\ge0}\sum_{j=1}^k
              \frac{(j)_m p_{k,j}(n)}{(z+n)^{j+m}}.
 \label{eq:all-derivatives}
\end{equation}
The series is normally convergent.  Its left side is explicitly
\begin{equation}
 \sum_{m_1+\cdots+m_k=m}
          \frac{m!}{m_1!\cdots m_k!}
                         \prod_{i=1}^k\psi^{(m_i)}(z),
 \label{eq:polygamma-multinomial}
\end{equation}
where $\psi^{(0)}=\psi$ and $(j)_m$ is the rising factorial.
\end{corollary}

\begin{proof}
Differentiate \eqref{eq:ML} normally on compact subsets.  The constant
subtractions disappear after the first derivative and each pole gives
the displayed rising factorial.  The denominator exponent is at least
two, so \eqref{eq:p-growth} proves normal convergence directly as well.
The multinomial identity is the ordinary product rule.
\end{proof}

For example, with $b_n=H_n-\gamma$ the case $k=2,m=1$ gives
\[
 \psi(z)\psi'(z)
     =-\sum_{n\ge0}\left\{\frac1{(z+n)^3}
                             -\frac{b_n}{(z+n)^2}\right\}.
\]
This agrees with the classical polygamma pole series at $k=1$ and
extends it to polynomial products without an undetermined analytic
term.  The full formulas apply at arbitrary complex $z$ in their
specified domains, not only at integer arguments.

\section{Ordered depth generators and exact cancellations}
\label{sec:ordered}

The preceding spectral derivatives retain information about the order of
the summation indices.  A complementary question is whether selected
combinations over \emph{all depths} admit finite evaluations.  We give
such evaluations while keeping one tagged exponent arbitrary.  The
resulting formulas connect ordered harmonic sums to ordinary Hurwitz-zeta
derivatives, generalized Stieltjes constants, and $\log\Gamma$.

The entire interpolation used here is established: its one-endpoint
form is the interpolant of Ihara, Nakamura, and Yamamoto
\cite{INY}, and its relative formulation already appears in the
incoming independent-order report.  Gamma formulas for depth-generating
series have substantial antecedents, including Young's work
\cite{Young}.  The purpose of this section is to derive a common
two-marker formula, prove its convergence and continuation, and extract
explicit cancellations useful for the present research program.

\subsection{Relative ordered harmonic sums}

For $c>0$ and a word $w=(s_1,\ldots,s_d)$, write
\[
 Z_c(w)=\sum_{n_1>\cdots>n_d\ge0}
              \prod_{j=1}^d(n_j+c)^{-s_j},
 \qquad Z_c(\varnothing)=1.
\]
All initial manipulations can be made in the smaller absolute-convergence
domain $\Re s_j>1$.  Let $F_c$ be the formal word series with these
coefficients and constant coefficient one.  Products concatenate words:
the coefficient of $w$ in $FG$ is $\sum_{w=uv}F(u)G(v)$.  Define
\begin{equation}\label{ord:relative}
             \Hrel_{a,b}=F_a^{-1}F_b,\qquad
             h_{a,b}(s)=\zeta(s,b)-\zeta(s,a).
\end{equation}
This specifies every coefficient by a finite algebraic expression.  For
example,
\[
 \Hrel_{a,b}(s)=h_{a,b}(s),\qquad
 \Hrel_{a,b}(s,t)=Z_b(s,t)-Z_a(s,t)-\zeta(s,a)h_{a,b}(t).
\]
The identity $\Hrel_{a,1}(w)=\Psi(w;a-1)$ identifies the one-endpoint
coefficients with \cite[Proposition~2.1]{INY}; their joint entireness in
the spectral indices is \cite[Corollary~4.2]{INY}.  It also proves the
relative assertion without a new regularization prescription, because
\[
                 \Hrel_{a,b}=\Hrel_{a,1}\Hrel_{b,1}^{-1}
\]
and each coefficient of the inverse is a finite polynomial in the
coefficients of $\Hrel_{b,1}$.  We always use these entire spectral continuations.
When $a=b+N$, $N\in\mathbb Z_{\ge0}$, cancellation of the ordered products
gives
\begin{equation}\label{ord:finite}
 \Hrel_{b+N,b}(s_1,\ldots,s_d)
   =\sum_{N>n_1>\cdots>n_d\ge0}
             \prod_{j=1}^d(n_j+b)^{-s_j}.
\end{equation}
Thus this is an interpolation of strict finite harmonic sums, with the
largest index on the left.  Write $1^p$ for a block of $p$ copies of $1$,
including the empty block when $p=0$.

\subsection{A two-marker theorem}

\begin{theorem}[One tagged exponent and two depth markers]\label{ord:two}
Let $a,b>0$, and put $\rho=\min(a,b)$.  The series
\begin{equation}\label{ord:Gdef}
 \mathscr G_s^{a,b}(x,y)
   =\sum_{p,q\ge0}x^py^q\Hrel_{a,b}(1^p,s,1^q)
\end{equation}
converges normally on compact subsets of
$\mathbb C_s\times\{|x|<\rho,\ |y|<\rho\}$.  Define, initially when
$\Re(s+x-y)>1$,
\begin{equation}\label{ord:kernel}
 S_s(c;x,y)=
   \sum_{n\ge0}(n+c)^{1-s}
       \frac{\Gamma(n+c+y)}{\Gamma(n+c+1+x)}.
\end{equation}
Then
\begin{equation}\label{ord:twoformula}
 \boxed{\;
 \mathscr G_s^{a,b}(x,y)
 =\frac{\Gamma(a+x)\Gamma(b)}
        {\Gamma(a)\Gamma(b+y)}
       \bigl(S_s(b;x,y)-S_s(a;x,y)\bigr).
 \;}
\end{equation}
The difference on the right has an entire continuation in $s$, joint
with the marking variables in the stated bidisc.  The completed formula
also has a meromorphic continuation in $x,y$.  On $y=0$ the larger domain
$|x|<a$, with arbitrary $b>0$, is a domain of normal convergence of the
one-marker series.  All spectral derivatives may be taken term by term
on these domains.
\end{theorem}

\begin{proof}
For the finite interval \eqref{ord:finite}, choosing the index $n+b$ of
the tagged exponent and then the subsets above and below it gives
\begin{align}
 \mathscr G_s^{b+N,b}(x,y)
 &=\sum_{n=0}^{N-1}(n+b)^{-s}
       \prod_{k=n+1}^{N-1}\left(1+\frac{x}{k+b}\right)
       \prod_{k=0}^{n-1}\left(1+\frac{y}{k+b}\right) \notag\\
 &=\frac{\Gamma(b+N+x)\Gamma(b)}
          {\Gamma(b+N)\Gamma(b+y)}
    \sum_{n=0}^{N-1}(n+b)^{1-s}
          \frac{\Gamma(n+b+y)}{\Gamma(n+b+1+x)}.
 \label{ord:finitegamma}
\end{align}
Removable Gamma specializations are interpreted through the finite
product in the first line.

Agreement on these finite intervals alone would not justify arbitrary
endpoints.  For that step use triangular matrices
\[
 E_c(n)=
 \begin{pmatrix}
  1+x/(n+c)&(n+c)^{-s}\\
  0&1+y/(n+c)
 \end{pmatrix},
 \qquad P_{c,M}=E_c(M-1)\cdots E_c(0).
\]
Their upper-right entry records exactly one tagged letter, with $x$
counting the letters before it and $y$ those after it.  Consequently
the upper-right entry of $P_{a,M}^{-1}P_{b,M}$ records the corresponding
relative coefficients.

To pass rigorously to exponent $1$, first replace every background
factor $(n+c)^{-1}$ by $(n+c)^{-\tau}$, with real $\tau>1$.
The infinite products then give the coefficients
$\Hrel_{a,b}(\tau^p,s,\tau^q)$ in an absolute-convergence region.
Fix $0<\delta<\rho/2$ and $|x|,|y|<\delta$.  For $\tau\ge1$
sufficiently close to $1$, the logarithms of the diagonal
\emph{quotient} products have differences $O(n^{-\tau-1})$.
After division by the first diagonal entry, each off-diagonal summand
is bounded by
$C(n+1)^{-\Re s+2\delta}$, locally uniformly in the parameters.
Indeed its product factors have logarithmic absolute value at most
$(|x|+|y|)\log(n+1)+O(1)$.
Thus the products and off-diagonal sums converge uniformly for
$\Re s>1+2\delta$ as $\tau\downarrow1$.  Cauchy coefficient extraction
commutes with this limit.  Entireness of the relative coefficients
therefore identifies its coefficients with
$\Hrel_{a,b}(1^p,s,1^q)$.

For completeness, the first and second diagonal entries of $P_{c,M}$
are $A_c(M,x)$ and $A_c(M,y)$, where
\[
 A_c(M,z)=
   \frac{\Gamma(M+c+z)\Gamma(c)}
        {\Gamma(M+c)\Gamma(c+z)}.
\]
Its upper-right entry is
\[
 \frac{\Gamma(M+c+x)\Gamma(c)}
      {\Gamma(M+c)\Gamma(c+y)}
 \sum_{n=0}^{M-1}(n+c)^{1-s}
       \frac{\Gamma(n+c+y)}{\Gamma(n+c+1+x)}.
\]
Inserting these expressions in the inverse of a triangular matrix and
using $\Gamma(M+c+z)/\Gamma(M+c)\sim M^z$ proves
\eqref{ord:twoformula} on the initial region just specified.

It remains to justify the claimed continuation and the convergence
of its Taylor series.  The Gamma-ratio expansion
\cite[\S5.11(iii)]{DLMF} is locally uniform in the markers:
\[
 \frac{\Gamma(t+y)}{\Gamma(t+1+x)}
  =t^{y-x-1}
    \left(\sum_{j=0}^{J-1}c_j(x,y)t^{-j}+O(t^{-J})\right),
 \qquad c_0=1.
\]
The coefficients do not depend on the endpoint; for instance,
$c_1=(y-x-1)(y+x)/2$.  Let $r_J(t)$ be the remainder after multiplying
this expansion by $t^{1-s}$.  Subtraction of the same asymptotic terms
at both endpoints yields
\begin{align}
 S_s(b;x,y)-S_s(a;x,y)
  ={}&\sum_{j=0}^{J-1}c_j(x,y)h_{a,b}(s+x-y+j)\notag\\
    &+\sum_{n\ge0}\bigl(r_J(n+b)-r_J(n+a)\bigr).
 \label{ord:subtract}
\end{align}
The remainder series is normally convergent whenever
$\Re(s+x-y)+J>1$.  Each $h_{a,b}$ is entire: the two simple Hurwitz-zeta
poles have the same residue and cancel.  Increasing $J$ therefore
continues the difference to every $s$, jointly holomorphically on
compact subsets of the stated bidisc.  The finitely many small-$n$
terms present no singularity there.

This continuation has Taylor coefficients in $x,y$ that are entire in
$s$.  They agree with those in \eqref{ord:Gdef} on the initial
half-plane, hence agree everywhere.  The Taylor series of a holomorphic
function on a bidisc converges normally on its compact subsets, proving
the series assertion.  The same subtraction, away from the discrete
Gamma-pole sets, gives meromorphic continuation in the markers.
Finally, at $y=0$ the kernel
$\Gamma(t)/\Gamma(t+1+x)$ is entire in $x$.  The only possible marking
poles of the completed expression then come from $\Gamma(a+x)$.
This proves holomorphy and Taylor convergence for $|x|<a$.
\end{proof}

\subsection{Integer gaps and ordinary zeta derivatives}

The gamma quotient in \eqref{ord:kernel} becomes a polynomial when
the markers have a suitable integer difference.  To state every such
case uniformly, define $c_{m,k}(x)$ by
\[
          (t+x+1)_m=\sum_{k=0}^m c_{m,k}(x)t^k,
 \qquad m\in\mathbb Z_{\ge0},
\]
where $(v)_m=v(v+1)\cdots(v+m-1)$ and $(v)_0=1$.

\begin{corollary}[Finite reduction at an integer marker gap]
\label{ord:gap}
For every complex $s$,
\begin{equation}\label{ord:gapformula}
 \mathscr G_s^{a,b}(x,x+m+1)=
 \frac{\Gamma(a+x)\Gamma(b)}
      {\Gamma(a)\Gamma(b+x+m+1)}
       \sum_{k=0}^m c_{m,k}(x)h_{a,b}(s-1-k).
\end{equation}
This is an ordinary convergent depth sum when the markers belong to the
convergence domain of Theorem~\ref{ord:two}; elsewhere it is an identity
of the specified meromorphic continuations.
In particular, for an integer $m\ge1$, $a>m$, and $b>0$,
\begin{equation}\label{ord:negative}
 \boxed{\;
 \sum_{p\ge0}(-m)^p\Hrel_{a,b}(1^p,s)
 =\frac{\Gamma(a-m)}{\Gamma(a)}
    \sum_{k=0}^{m-1}(-1)^{m-1-k}
           \genfrac{[}{]}{0pt}{}{m}{k+1}h_{a,b}(s-1-k).
 \;}
\end{equation}
Here $\genfrac{[}{]}{0pt}{}{m}{j}$ is an unsigned Stirling number of the
first kind.  The series in \eqref{ord:negative} converges normally
locally in $s$.
\end{corollary}

\begin{proof}
On $y=x+m+1$ the kernel quotient is exactly $(t+x+1)_m$.
Equation \eqref{ord:gapformula} follows first where the defining
Dirichlet series converge, and then everywhere by
\eqref{ord:subtract}.  For \eqref{ord:negative}, set $x=-m$, $y=0$
and use
\[
 \frac{\Gamma(t)}{\Gamma(t+1-m)}
   =\prod_{j=1}^{m-1}(t-j)
   =\sum_{k=0}^{m-1}(-1)^{m-1-k}
                      \genfrac{[}{]}{0pt}{}{m}{k+1}t^k.
\]
The condition $a>m$ places this specialization strictly inside the
one-marker Taylor disc.
\end{proof}

The first two reductions, including their normal-convergence conditions,
are particularly simple:
\begin{align}
 \sum_{p\ge0}(-1)^p\Hrel_{a,b}(1^p,s)
    &=\frac{h_{a,b}(s-1)}{a-1}, && a>1, \label{ord:mone}\\
 \sum_{p\ge0}(-2)^p\Hrel_{a,b}(1^p,s)
    &=\frac{h_{a,b}(s-2)-h_{a,b}(s-1)}
            {(a-2)(a-1)}, && a>2. \label{ord:mtwo}
\end{align}
All their spectral derivatives are finite combinations of ordinary
Hurwitz-zeta derivatives.  At a canceled pole they become Stieltjes
differences.  More explicitly, use the convention
\[
 \zeta(1+u,c)=\frac1u+
       \sum_{k\ge0}\frac{(-1)^k}{k!}\gamma_k(c)u^k.
\]
Differentiating \eqref{ord:mone} at $s=2$ gives, for $a>1$, $b>0$,
\begin{equation}\label{ord:stieltjes}
 \boxed{\;
 \sum_{p\ge0}(-1)^p
    \left.\frac{\partial^k}{\partial s^k}
          \Hrel_{a,b}(1^p,s)\right|_{s=2}
  =\frac{(-1)^k\bigl(\gamma_k(b)-\gamma_k(a)\bigr)}{a-1}.
 \;}
\end{equation}
For $k=0$ this is $(\psi(a)-\psi(b))/(a-1)$.
At the adjacent point $s=1$, the classical relation
$\zeta'(0,c)=\log\Gamma(c)-\tfrac12\log(2\pi)$
\cite[\S25.11]{DLMF} instead yields
\begin{equation}\label{ord:gamma}
 \boxed{\;
 \sum_{p\ge0}(-1)^p
    \left.\frac{\partial}{\partial s}
          \Hrel_{a,b}(1^p,s)\right|_{s=1}
       =\frac{\log\Gamma(b)-\log\Gamma(a)}{a-1}.
 \;}
\end{equation}
For $a=N+1$, $b=1$, $N\ge1$, the sum is finite and its value is
$-\log(N!)/N$.  These statements evaluate ordered combinations over
depth.  They do not remove the nonsymmetric spectral information in
each separate fixed-depth coefficient.

\subsection{A finite-interval formula with several tagged exponents}

\begin{corollary}[A finite staircase of markers]\label{ord:staircase}
Let $a=b+N$, $b>0$, $N\in\mathbb Z_{\ge0}$, and let $r\ge1$ with
$a>r$.  For arbitrary complex $s_1,\ldots,s_r$,
\begin{equation}\label{ord:stairformula}
 \begin{split}
 &\sum_{k_0,\ldots,k_{r-1}\ge0}
       \left(\prod_{j=0}^{r-1}(j-r)^{k_j}\right)
       \Hrel_{a,b}(1^{k_0},s_1,1^{k_1},s_2,\ldots,
                                  1^{k_{r-1}},s_r)\\
 &\hspace{35mm}
       =\frac{\Gamma(a-r)}{\Gamma(a)}
                    \Hrel_{a,b}(s_1-1,\ldots,s_r-1).
 \end{split}
\end{equation}
The displayed sum is finite: all terms of depth greater than $N$ vanish.
\end{corollary}

\begin{proof}
Introduce independent markers $x_0,\ldots,x_r$ for the blocks of ones
before, between, and after the tagged exponents.  The finite-product
calculation in \eqref{ord:finitegamma}, now made at each selected
tagged index $t_j=n_j+b$, gives their joint generator as
\[
 \frac{\Gamma(a+x_0)\Gamma(b)}
      {\Gamma(a)\Gamma(b+x_r)}
 \sum_{N>n_1>\cdots>n_r\ge0}
       \prod_{j=1}^r
           t_j^{1-s_j}
           \frac{\Gamma(t_j+x_j)}
                {\Gamma(t_j+1+x_{j-1})}.
\]
Set $x_j=j-r$.  Every quotient inside the sum becomes one, the last
marker is zero, and the prefactor becomes $\Gamma(a-r)/\Gamma(a)$.
The calculation may first be made with $b>r$, so that no intermediate
Gamma factor is singular.  Both sides are real analytic on the
connected interval $b>\max(0,r-N)$, so the identity extends throughout
that interval.  Equivalently one may use the finite products to fill
all removable intermediate specializations.
\end{proof}

For example, with the same finite-interval hypothesis and $a>2$,
\[
 \sum_{p,q\ge0}(-2)^p(-1)^q
          \Hrel_{a,b}(1^p,s,1^q,t)
       =\frac{\Hrel_{a,b}(s-1,t-1)}{(a-1)(a-2)}.
\]
At $s=t=2$ its value is
\[
 \frac{(\psi(a)-\psi(b))^2-\zeta(2,b)+\zeta(2,a)}
           {2(a-1)(a-2)}.
\]
Corollary~\ref{ord:staircase} is asserted for finite intervals only.
An extension to independent endpoints with several unbounded depth
markers would require its own convergence argument.

\subsection{A polylogarithmic companion}

Use the one-variable multiple polylogarithm convention
\[
 \Li_{s_1,\ldots,s_d}(z)
    =\sum_{n_1>\cdots>n_d\ge1}
        \frac{z^{n_1}}{n_1^{s_1}\cdots n_d^{s_d}},
 \qquad |z|<1.
\]
This convention places the argument on the largest index, as in the
ordered harmonic sums above.

\begin{proposition}[A depth generator and finite-depth evaluations]
\label{ord:polylog}
For $|z|<1$ and arbitrary $s,u\in\mathbb C$, the series
\[
                  \mathscr F_s(z;u)
                    =\sum_{p\ge0}u^p\Li_{1^p,s}(z)
\]
converges normally on compact subsets of these variables and satisfies
\begin{equation}\label{ord:polygen}
 \boxed{\;
 \mathscr F_s(z;u)
   =(1-z)^{-u}
           \int_0^z\frac{(1-t)^u\Li_{s-1}(t)}{t}\,\dd t.
 \;}
\end{equation}
The branch of $\log(1-z)$ vanishes at zero.  In particular,
\begin{align}
 \sum_{p\ge0}(-1)^p\Li_{1^p,s}(z)
   &=(1-z)\bigl(\Li_s(z)+\Li_{1,s-1}(z)\bigr),
       \label{ord:polyone}\\
 \sum_{p\ge0}(-2)^p\Li_{1^p,s}(z)
   &=(1-z)^2\bigl(\Li_s(z)+\Li_{1,s-1}(z)
                                  -\Li_{1,s-2}(z)\bigr)\notag\\
   &\hspace{8mm}+z(1-z)\Li_{s-1}(z).
       \label{ord:polytwo}
\end{align}
More generally every negative integer $u=-m$ admits a finite expression
using ordinary polylogarithms and $\Li_{1,s-j}(z)$ with rational
coefficients in $z$.  All spectral derivatives of these identities
are valid on the same disc.
\end{proposition}

\begin{proof}
For $|u|\le R$ and $s$ in a fixed compact set, the absolute sum of the
coefficients attached to $z^n$ is bounded by
$n^C\prod_{k=1}^n(1+R/k)=O(n^{C+R})$ for a suitable $C$.
Indeed one may first choose the tagged index and then independently
choose the larger background indices.  Summation against $|z|^n$
proves normal convergence, including in $u$ and $s$.

Differentiation of the defining series gives
\[
 \frac{\partial\mathscr F_s}{\partial z}
    =\frac{\Li_{s-1}(z)}z+\frac{u}{1-z}\mathscr F_s(z;u),
 \qquad \mathscr F_s(0;u)=0.
\]
Multiplication by the integrating factor $(1-z)^u$ proves
\eqref{ord:polygen}.  At $u=-m$ use
\[
 \frac1{t(1-t)^m}=\frac1t+\sum_{j=1}^m\frac1{(1-t)^j}.
\]
Thus the integral is $\Li_s(z)+\sum_{j=1}^m I_j(s-1,z)$, where
$I_j(v,z)=\int_0^z\Li_v(t)(1-t)^{-j}\,\dd t$.
Here $I_1(v,z)=\Li_{1,v}(z)$, and integration by parts gives the finite
recursion, for $j\ge2$,
\[
 I_j(v,z)=\frac{\Li_v(z)}{(j-1)(1-z)^{j-1}}
       -\frac{\Li_v(z)+\sum_{k=1}^{j-1}I_k(v-1,z)}{j-1}.
\]
This proves the general finite-depth assertion and, at $m=1,2$,
the two displayed formulas.  Normal convergence permits all derivatives
in the tagged exponent.
\end{proof}

These polylogarithmic formulas, like \eqref{ord:negative}, reduce a
specified depth-generating combination.  They give exact primitives
and spectral derivatives with controlled depth, while leaving open
the separate reduction problems for individual ordered coefficients.

\section{Source audit and precise corrections}
\label{powers:sec:audit}

\subsection{Scope of the repository audit}

The source baseline is ProveIt commit
\texttt{b7667ae40a2957eed9a53d8559d53a3ee20dd1c1}, including the
current manuscript and the incoming archives inspected for this
continuation \cite{Repo}. The review concentrated on the polynomially
weighted resolvent, the multiple-Gamma normalization, continuous
resolvent powers, collision finite parts, and the quartic
power-to-harmonic bridge \cite{Quartic,Barnes}. No new false theorem
was identified in the selected repository statements used here.
This is a targeted mathematical audit, not a certification of every
chapter of the complete manuscript. The two corrections below concern
precisely identified external sources.

The finite-part convention remains the original unit coordinate:
\(\FP\int_0^1 f(x)\,\dd x\) retains the constant coefficient of
\(\int_\epsilon^1 f(x)\,\dd x\) after its explicit local power and
logarithmic terms are separated. It is not an unspecified analytic
regularization. Likewise, a constant coefficient in a spectral Laurent
series or in a polynomial asymptotic expansion in \(\log N\) is a
different operation until a comparison theorem proves equality. In the
present argument the normalized contour formula and finite summation
by parts provide that comparison. Matching principal parts alone would
leave an undetermined entire remainder and would not prove the moment
identity.

\subsection{A reversed pair of integrals in a Tornheim source}

The author-hosted version of Bailey and Borwein's
\emph{Computation and theory of Mordell--Tornheim--Witten sums II},
dated 5 May 2014, contains a sign reversal in equation~(110), on printed
page~24 \cite{BB2014}. Its preceding equation~(109) is correct. The
correction is recorded specifically for this inspected dated version;
no claim is made here about an independently uninspected final
typesetting.

For \(\Re a>0\), put
\begin{align}
 A(a)&=\frac1{\Gamma(a+1)}\int_0^1
       \frac{(-\log u)^a[-\log(1-u)]}{1-u}\,\dd u,
       \label{powers:audit:A}\\
 B(a)&=\int_0^1\Li_a(u)
                  \frac{-\log(1-u)}u\,\dd u.
       \label{powers:audit:B}
\end{align}
The power of \(-\log u>0\) uses its real logarithm. Both integrals
converge absolutely and depend holomorphically on \(a\) in this
half-plane.

\begin{proposition}[Corrected integral identity]
\label{powers:audit:BB-correction}
For \(\Re a>0\),
\begin{equation}
 \boxed{\displaystyle
 \zeta(a+2)=B(a)-A(a).}
 \label{powers:audit:correct-order}
\end{equation}
Equation~(110) of the specified source prints \(A(a)-B(a)\).
\end{proposition}

\begin{proof}
Write \(H_m=\sum_{j=1}^m j^{-1}\) and
\(T(a,1,1)=\sum_{m,n\ge1}m^{-a}n^{-1}(m+n)^{-1}\).
For each positive integer \(m\), the exact finite calculation
\[
 \sum_{n=1}^N\frac1{n(m+n)}
 =\frac{H_N+H_m-H_{N+m}}m
\]
gives \(\sum_{n\ge1}1/[n(m+n)]=H_m/m\). Therefore
\begin{equation}
 T(a,1,1)=\sum_{m\ge1}\frac{H_m}{m^{a+1}}
         =\zeta(a+1,1)+\zeta(a+2),
 \label{powers:audit:direct-T}
\end{equation}
where \(\zeta(s,1)=\sum_{m\ge1}H_{m-1}m^{-s}\).
These series converge absolutely when \(\Re a>0\).
Expanding the polylogarithm and logarithm in
\eqref{powers:audit:B} gives \(B(a)=T(a,1,1)\), while
\[
 \frac{-\log(1-u)}{1-u}=\sum_{n\ge1}H_nu^n,
 \qquad
 \int_0^1u^n(-\log u)^a\,\dd u
       =\frac{\Gamma(a+1)}{(n+1)^{a+1}}
\]
give \(A(a)=\zeta(a+1,1)\). Absolute convergence justifies all
interchanges, proving \eqref{powers:audit:correct-order}.
\end{proof}

There is an exact elementary counterexample to the printed sign.
At \(a=1\), substitution followed by the geometric series gives
\[
 A(1)=\int_0^1\frac{\log v\log(1-v)}v\,\dd v=\zeta(3),
 \qquad
 B(1)=\int_0^1\frac{\log^2v}{1-v}\,\dd v=2\zeta(3).
\]
Thus its printed right side is \(-\zeta(3)\), whereas its left
side is \(\zeta(3)>0\). This corrects the integral line; it does
not invalidate equation~(109) or the spectral derivative identities
obtained from that correct equation. The proof above also establishes
the whole half-plane identity directly, without a Carlson-theorem
argument.

\subsection{A distinction between multiple and ordinary zeta reductions}

On printed page~205 of \emph{Nonlinear Euler sums}, Mez\H{o}
states that it is not known whether all the nonlinear sums under
discussion can be expressed as multiple zeta values \cite{Mezo2014}.
For convergent products of positive-integer generalized harmonic
numbers, this target class is too broad: a finite multiple-zeta
reduction follows directly from ordering their summation indices.

To make the correction precise, use the strict finite sums
\[
 H_n(b_1,\ldots,b_\ell)
 =\sum_{n\ge j_1>\cdots>j_\ell\ge1}
       j_1^{-b_1}\cdots j_\ell^{-b_\ell}.
\]
For positive integers \(a_1,\ldots,a_d\), expand
\(\prod_{i=1}^d H_n^{(a_i)}\) by the ordered blocks of equal
indices. Every ordered set partition
\((B_1,\ldots,B_\ell)\) of \(\{1,\ldots,d\}\) contributes
\(H_n(b_1,\ldots,b_\ell)\), with
\(b_j=\sum_{i\in B_j}a_i\). This is the finite stuffle product.
For \(q\ge2\), each resulting term satisfies
\begin{equation}
 \sum_{n\ge1}\frac{H_n(b_1,\ldots,b_\ell)}{n^q}
 =\zeta(q,b_1,\ldots,b_\ell)
       +\zeta(q+b_1,b_2,\ldots,b_\ell).
 \label{powers:audit:finite-stuffle}
\end{equation}
The two terms distinguish \(n>j_1\) and \(n=j_1\).
All sums are absolutely convergent; hence this is a finite exact
reduction for every such nonlinear Euler sum.

For example,
\begin{align}
 \sum_{n\ge1}\frac{H_n^2}{n^2}
 &=2\zeta(2,1,1)+2\zeta(3,1)+\zeta(2,2)+\zeta(4)
   =\frac{17}{4}\zeta(4).
 \label{powers:audit:example}
\end{align}
The first equality already demonstrates the required finite
multiple-zeta reduction; the last equality uses the familiar
weight-four evaluations. A precise replacement for the source's
wording should identify a smaller target, such as ordinary zeta
values and their products, or a specified class of linear Euler
sums. Formula~\eqref{powers:audit:finite-stuffle} does not establish
either stronger reduction in general. It is a classical algebraic
mechanism, not a new arithmetic-independence result.

\subsection{Unchanged conjecture status}

The fifth-power identity and the algorithm for arbitrary integer
powers establish the finite harmonic spectral formulas stated in
this article. They do not prove that every retained spectral
derivative reduces to ordinary constants. A nonzero coordinate in a
specified formal presentation also does not establish arithmetic
independence of its numerical value. The current Gaussian
\(S_6\) and revised \(S_8\) identities in the pinned manuscript
remain conjectures \cite{Repo}; none of the arguments here changes
their status.

\section{Verification and reproducibility}
\label{sec:verification}

The proofs above supply the analytic justification.  Computation was
used to check finite algebra, branch and indexing conventions, and
numerical evaluations by independent routes.  It was not used to
replace the contour estimates, continuation arguments, or the
cutoff-to-Laurent comparison.

\subsection{Exact finite algebra}

The accompanying \texttt{code/power\_algebra.py} implements the coefficient
construction and strict-word stuffle recursion with exact rational
symbolic coefficients.  The script \texttt{code/exact\_checks.py}
verifies the finite difference recurrence through power eight, and
compares the harmonic-word expansion against actual finite nested
sums at several distinct indices.  It also checks the full fifth
spectral support, the low-power scalar reductions, the closed residue
polynomial, and the canonical fifth polynomial.  Its output is stored
in \texttt{data/exact\_checks.json}; the reusable word tables are in
\texttt{data/power\_coefficients.json}.

These finite checks have a precise limitation.  For example, verifying
the recurrence at finitely many powers does not prove its assertion
for arbitrary $k$; Proposition~\ref{prop:difference} does.  Likewise,
symbolic reduction of a finite polynomial does not justify exchanging
an infinite sum with an endpoint finite part.  That step is supplied
by Theorems~\ref{thm:ML}, \ref{thm:contour}, and \ref{thm:all-power}.

\subsection{An independent fifth-power calculation}

The script \texttt{code/independent\_fifth.py} imports none of the main
coefficient implementation.  It reconstructs the fifth pole data,
the higher-pole contribution, and the cutoff constant through its own
finite stuffle calculation.  It then evaluates $Q_5$ in two ways.
The first uses local Taylor subtraction of the pole of $\psi(x)^5$,
so no singular numerical quadrature is performed without subtraction.
The second uses the ordinary arithmetic series
\eqref{eq:Q5ordinary}, with an independently generated Bernoulli and
Hurwitz-zeta tail.  With 90 working decimal digits, cutoff $N=96$,
and tail powers through $n^{-36}$, the two routes gave
\begin{equation}
 Q_5\approx
 -8.69030396262453826629558966236565357677029354492388429266.
 \label{eq:q5decimal}
\end{equation}
Their absolute discrepancy was $9.60\times10^{-59}$.
These are floating-point diagnostics, not an interval-certified
enclosure.  The decimal in \eqref{eq:q5decimal} is therefore reported
as a diagnostic approximation; the exact assertion is
Theorem~\ref{thm:q5}.

The ordered-generator checks compare the Gamma formulas with direct
finite nested sums, including nonreal tagged exponents and nonintegral
lower endpoints.  They also check Stieltjes and log-Gamma derivative
consequences, the two-tag staircase, and the polylogarithmic formulas
at negative markers $-1$ and $-2$.  All 23 numerical comparisons in
\texttt{code/verify\_ordered\_generators.py} passed at 70 working
digits; the largest absolute residual was $4.92\times10^{-70}$.
Their defining finite sums and power-series evaluations are separate
from the closed formulas being checked.  These are again
non-interval diagnostics.

\subsection{How to reproduce the package}

The standalone \LaTeX{} file builds with an ordinary \texttt{pdflatex}
installation and the listed standard packages; no external images,
bibliography service, or shell escape is required.  The modular source
fragments are included to facilitate repository integration, while the
single-file source is sufficient for rebuilding the PDF.  Python
requirements are \texttt{sympy} and \texttt{mpmath}; the versions used
for the recorded checks are given in their JSON receipts.  Run
\begin{verbatim}
python3 code/exact_checks.py
python3 code/independent_fifth.py
python3 code/verify_ordered_generators.py
python3 code/build_article.py
\end{verbatim}
from the package directory.  The optional Gaussian relation experiment
has separate instructions and requires a C++17 compiler.

The integration directory contains a claim ledger, the pinned source
revision and hashes, proposed placement notes, and the exact scope of
the corrections.  It deliberately distinguishes proved additions,
recovered antecedents, diagnostics, and unresolved targets.  Nothing
in the package changes the conjecture status of unrelated source
statements.

\section{Further exact research questions}
\label{sec:questions}

The strongest next questions concern finite identities and arithmetic
reduction.  Estimates are useful to justify operations, but are not
the principal intended deliverable.

\subsection{The first mixed spectral coordinate}

Determine whether $\Omega_4$ in \eqref{eq:mixed-def} can be expressed
in a specified algebra generated by ordinary zeta derivatives,
ordinary Stieltjes constants, and explicitly named Gamma or Barnes
periods.  Equations~\eqref{eq:OmegaT}--\eqref{eq:Tpolylog} give
three concrete starting points: nonsymmetric derivatives of double
zeta functions, a Tornheim derivative, and a dilogarithm integral.
An answer should identify which independent-looking coordinates are
actually removed.  Merely renaming the dilogarithm integral does not
constitute an arithmetic reduction; a failed integer-relation search
does not prove nonreducibility.

One tractable first task is to study identities for
$\partial_1Z(2,2)$ and $(\partial_1-\partial_2)Z(3,1)$ together,
instead of trying to evaluate $D_{1,3}$ in isolation.  Differentiated
stuffle determines symmetric combinations, whereas the second
expression is explicitly antisymmetric in the spectral directions.
A new identity must address that remaining direction.

\subsection{Scalar polynomials and the growth of spectral support}

The normalized fifth polynomial in
Corollary~\ref{powers:cor:canonical-fifth} cancels the displayed
lower height-one coefficients, but retains $\Omega_4$.  For general
$k$, determine the image of scalar polynomials
$P(\psi)$ under the finite-word spectral map of
Theorem~\ref{thm:algorithm}.  One can first ask a wholly finite
algebraic question: in a stated presentation modulo specified stuffle
and differential relations, which coordinates can be eliminated by
lower powers?

A useful all-order result would give a recurrence for normalized
monic polynomials and a proved description of their remaining word
support.  The stronger assertion that only the single highest
height-one Laurent coefficient remains is not established by the
degree-four cancellation and is not established here at degree five.
An arithmetic simplification of the mixed coordinate could change
the conclusion; formal coordinates must not be silently promoted to
independent numbers.

\subsection{Explicit sixth and seventh moments}

The included coefficient tables through power eight make the next
cases finite, but a table is not yet a well-organized theorem.
Carry out $Q_6$ and $Q_7$ with a reduced, explicitly chosen spectral
coordinate family.  The all-power formula already supplies the
boundary constants $C_6=-\pi^6/7$ and $C_7=0$, and the complete
Stieltjes subtraction is explicit.  The remaining task is to organize
the ordinary multiple-zeta values in $J_k$ and the word derivatives
in $D_k$ without losing diagonal terms.

Compare the resulting support with the fifth mixed coordinate.
Does a finite preferred family emerge at each fixed weight after
all differentiated stuffle identities are imposed?  The deliverable
should include exact coordinate changes and reproducible rational
certificates, not only a numerical basis.

\subsection{A common Hurwitz shift with correct endpoint terms}

For $a>0$, seek an analogue for
$\int_0^1\psi(a+x)^k\,\dd x$ in a common-shift multiple-Hurwitz
notation.  A derivation must retain both ordinary endpoints
$a$ and $a+1$.  Then study the singular limit $a\downarrow0$
using the original coordinate and compare it with $Q_k$.
Replacing $\gamma_j$ by $\gamma_j(a)$ in the unshifted formula
without deriving the new pole locations and endpoint contributions
does not answer this question.

The relative functions in Section~\ref{sec:ordered} suggest a useful
language for endpoint differences.  Their spectral indices are
entire, but that alone does not determine the finite part of a
shifted digamma product.  A precise theorem should specify the
shift in every denominator and prove the connection.

\subsection{Differential polynomials in digamma and polygamma}

Extend the method from $\psi^k$ to an arbitrary fixed differential
polynomial
$P(\psi,\psi',\ldots,\psi^{(r)})$ with constant coefficients.
Its local pole data are finite and its large-circle growth is
controlled, suggesting a generalized normalized expansion and a
finite-part reduction.  Derivatives may change the angular
constant substantially, so it should be computed rather than
assumed to be the same $C_k$.

A first exact target is the family
$\FP\int_0^1\psi(x)^a\psi'(x)^b\,\dd x$ with integers
$a,b\ge0$.  Integration by parts gives relations between these
moments, and the pole-coefficient method supplies their boundary
terms.  Classifying the independent moments after these identities
would connect the primitive formulas to a broader polygamma calculus.

\subsection{Several tags at independent endpoints}

Corollary~\ref{ord:staircase} is proved only for finite intervals.
For independent positive endpoints $a,b$, construct the full
upper-triangular matrix generator with several tagged exponents,
prove normal convergence in all depth markers, and determine
whether the integer staircase still collapses to a shifted ordered
relative sum.  A direct matrix and asymptotic subtraction argument
would be preferable to interpolation from integer intervals, which
does not by itself prove the extension.

If such a theorem holds, differentiating at several spectral poles
should produce exact combinations of multivariate harmonic Laurent
data.  The useful target is an explicitly evaluated mixed derivative
combination, with its continuation independent of the order of
taking removable limits.

\subsection{Boundary markers and removable Gamma singularities}

The one-marker theorem provides an ordinary depth sum at $x=-m$
when $a>m$.  Study the boundary $a=m$ and the exceptional points
$a=m-\ell$ through the actual generating function, separating
ordinary convergence, meromorphic continuation, and finite part.
At finite intervals the generating function is a polynomial and
some apparent Gamma poles are removable; for independent endpoints
this cancellation needs a proof.  Determine the finite limiting
identities and their log-Gamma or Stieltjes corrections in each
admissible case.

\subsection{Polylogarithmic cancellation at roots of unity}

The identities of Proposition~\ref{ord:polylog} hold normally on
$|z|<1$.  Analyze their boundary continuation at $z=1,-1,i,-i$,
with explicit conditions on the tagged exponent.  At $z=1$ the
prefactors and divergent depth sums interact, so a boundary identity
must be derived with a specified limiting operation.  At other roots
of unity, seek exact colored multiple-zeta formulas and their
spectral derivatives.  A rigorous boundary theorem may supply
relations useful to the Gaussian part of the manuscript.

\subsection{Gaussian relations beyond the searched family}

The computation in Appendix~\ref{app:s6} shows only that one
explicitly generated finite family does not close the current
$S_6$ target modulo one prime.  A productive next step is to derive
new exact functional equations or construct the full relevant
Cayley ideal at a higher filtration, then seek a rational
certificate for the target.  Increasing a numerical precision or
repeating the same finite-field search does not supply a proof.
The revised $S_8$ target requires its own exact treatment and must
not inherit a proved label from a fitted expression.

\subsection{Formal verification of the analytic normalization}

The finite coefficient recurrence and the strict-word conversion
are suitable for small proof-assistant components.  The most
valuable analytic components would formalize the full-circle
digamma estimate, normal convergence of the subtracted
Mittag--Leffler series, and the cutoff-to-Laurent lemma.
Together they would make the all-power theorem reusable without
re-proving an endpoint regularization separately at each degree.
An exact algebra certificate alone should not be advertised as
formal verification of the analytic theorem.

\appendix
\section{A bounded investigation of the current Gaussian \texorpdfstring{$S_6$}{S6} target}
\label{app:s6}

For clarity, the searched target is recorded independently of the
proved digamma results.  In this appendix use the colored convention
\[
 \Li_{a,b}(z_1,z_2)=
       \sum_{n>m\ge1}\frac{z_1^nz_2^m}{n^am^b},\qquad
 g_{a,b}=\Im\Li_{a,b}(i,1),
\]
and $\beta(s)=\sum_{n\ge0}(-1)^n/(2n+1)^s$.
The current short residual in the pinned manuscript is
\begin{align}
 R_6={}&S_6+\frac{12334}{527}g_{6,1}
          +\frac{384}{31}g_{5,2}+\frac{2136}{527}g_{4,3}
 \notag\\
 &-\frac{3179}{5713920}\pi^7
     +\frac{801}{2108}\beta(4)\zeta(3)+2\beta(6)\log2,
 \label{eq:S6residual}\\
 S_6={}&\sum_{n\ge0}\frac{(-1)^nH_n}{(2n+1)^6}.
 \notag
\end{align}
The conjecture is $R_6=0$; no such equality is proved here.

We extended the explicitly generated weight-seven relation family
to maximum total input depth five.  After imaginary projection,
primitive integer normalization, and duplicate removal, the family
had the following size:
\begin{center}
\begin{tabular}{lr}
\toprule
Relation family & Rows\\
\midrule
Double-shuffle relations & $14\,908$\\
Distribution relations and specified product lifts & $4\,938$\\
Lifts of the proved $S_2/S_4$ relations & $128$\\
Balanced Cayley relations and specified lifts & $2\,322$\\
\midrule
Total & $22\,296$\\
\bottomrule
\end{tabular}
\end{center}
There were $8\,786$ occupied imaginary word coordinates.  A completed
sparse elimination over $\mathbb F_{2147483647}$ found rank
$8\,705$ and left a nonzero $16$-coordinate representative of the
target.  The family was exhausted in approximately $328$ seconds
on the machine used, before the stipulated $600$-second cap.
Timing is not a mathematical invariant; the field, row family,
rank, and nonzero reduced target are the reproducible diagnostic.

The imaginary evaluation of the exact integer target is $128588R_6$.
A separate rational-arithmetic script verifies this normalization,
including the conversion $\beta(7)=61\pi^7/184320$.
The family generator, field eliminator, complete-run receipts,
normalization check, and reproduction note are included in
\texttt{exploratory/S6\_depth5\_bounded\_search.zip}.
The raw matrix is regenerated by the included code.

This establishes nonmembership in the \emph{specified finite-field
row span}.  It does not establish rational-span nonmembership:
an integer matrix can lose rank modulo the selected prime.
It does not establish nonvanishing of $R_6$, and it does not prove
the conjecture.  Moreover, the finite balanced Cayley family here
is not claimed to generate the full Cayley shuffle ideal at depth
five.  This distinction prevents comparison with stronger exact
obstruction statements in the incoming material from being
overstated.  The retained log is a reproducible computation record,
not a separately replayed compact rational certificate.

\subsection{A useful local attribution}

The canonical $S_2$ section would benefit from a citation immediately
beside its classical formula
\begin{equation}
 S_2=2\Im\Li_3\!\left(\frac{1+i}{2}\right)
       -G\log2-\frac{\pi(\log2)^2}{16}-\frac{\pi^3}{64},
 \qquad G=\beta(2).
 \label{eq:S2attribution}
\end{equation}
Xu and Wang, Section~3.4, Example~3.13 of the inspected arXiv
version~3 \cite{XuWang}, give the identical evaluation after
$R_{1,\bar2}=4S_2$ is substituted.  The broader manuscript already
cites that work; the proposed change is a local attribution
improvement, not a correction to the value or its proof.

\begin{thebibliography}{99}
\addcontentsline{toc}{section}{References}
\raggedright

\bibitem{Repo}
V. Reshetnikov and contributors,
\emph{ProveIt}, source revision
\texttt{b7667ae40a2957eed9a53d8559d53a3ee20dd1c1},
inspected 11 October 2026.
Manuscript and incoming reports:
\url{https://github.com/VladimirReshetnikov/ProveIt/tree/b7667ae40a2957eed9a53d8559d53a3ee20dd1c1/Analysis/Polylogarithms/docs/manuscript}
and
\url{https://github.com/VladimirReshetnikov/ProveIt/tree/b7667ae40a2957eed9a53d8559d53a3ee20dd1c1/docs/incoming}.

\bibitem{Quartic}
ProveIt research continuation,
\emph{Quartic Digamma Bridges, Harmonic Tail Moments, and Complex
Resolvent Powers}, 11 October 2026.
Incoming archive
\texttt{ProveIt\_Quartic\_Tails\_Resolvents\_2026-10-11.zip}
at the revision in \cite{Repo}, especially the quartic bridge
and further-research Questions~1--4.

\bibitem{Barnes}
ProveIt research continuation,
\emph{Barnes Resolvents and Higher Harmonic Identities},
11 October 2026.
Incoming archive
\texttt{ProveIt\_Barnes\_Resolvents\_and\_Harmonic\_Identities\_2026-10-11.zip}
at the revision in \cite{Repo}, especially the quartic harmonic
section and its normalization.

\bibitem{Independent}
ProveIt research continuation,
\emph{Independent Orders and Exact Reductions in
Polylogarithm--Zeta Calculus}, 11 October 2026.
Incoming archive
\texttt{ProveIt\_Independent\_Orders\_2026-10-11.zip}
at the revision in \cite{Repo}, especially the relative
harmonic interpolation and exact Gaussian relation framework.

\bibitem{DLMF}
NIST Digital Library of Mathematical Functions.
Sections~5.5 (functional relations), 5.7 (series expansions),
5.11 (asymptotics and Gamma ratios), 25.11 (Hurwitz zeta),
25.12 (polylogarithms), and 25.16 (multiple zeta values).
\url{https://dlmf.nist.gov/}.

\bibitem{IKZ}
K. Ihara, M. Kaneko, and D. Zagier,
Derivation and double shuffle relations for multiple zeta values,
\emph{Compositio Mathematica} \textbf{142} (2006), 307--338.
\url{https://doi.org/10.1112/S0010437X0500182X}.

\bibitem{INY}
K. Ihara, Y. Nakamura, and S. Yamamoto,
Interpolant of truncated multiple zeta functions.
arXiv:2407.20509 (2024); published in
\emph{The Ramanujan Journal} \textbf{67} (2025), Article~1.
\url{https://arxiv.org/abs/2407.20509}.
\url{https://doi.org/10.1007/s11139-025-01045-2}.
Proposition~2.1 and Corollary~4.2 are cited using the arXiv text.

\bibitem{Young}
P. T. Young,
Series of height one multiple zeta functions,
\emph{Integers} \textbf{24} (2024), A43.
\url{https://math.colgate.edu/~integers/y43/y43.pdf}.
\url{https://doi.org/10.5281/zenodo.11221606}.

\bibitem{BB2014}
D. H. Bailey and J. M. Borwein,
\emph{Computation and theory of Mordell--Tornheim--Witten sums II}.
Author-hosted version dated 5 May 2014, inspected at
\url{https://www.davidhbailey.com/dhbpapers/mtw2.pdf}.
The correction in Section~\ref{powers:sec:audit} is specifically to
equation~(110), printed page~24, in this version.

\bibitem{Mezo2014}
I. Mez\H{o}, Nonlinear Euler sums,
\emph{Pacific Journal of Mathematics} \textbf{272} (2014), no.~1,
201--226.
\url{https://doi.org/10.2140/pjm.2014.272.201}.
\url{https://msp.org/pjm/2014/272-1/pjm-v272-n1-p09-s.pdf}.
The wording discussed in Section~\ref{powers:sec:audit} is on
printed page~205.

\bibitem{XuWang}
C. Xu and W. Wang,
\emph{Dirichlet type extensions of Euler sums}.
arXiv:2009.11704v3, 12 April 2022,
Section~3.4, Example~3.13.
\url{https://arxiv.org/html/2009.11704v3}.

\end{thebibliography}
\end{document}

